\documentclass[11pt,a4paper]{article}
\makeatletter\def\input@path{{sections/notes/}{sections/paper/}}\makeatother
% LuaLaTeX. The optional font file is only a local build convenience.
\usepackage{luatexja-fontspec}
\IfFileExists{../local-fonts.tex}{\input{../local-fonts.tex}}{}
\usepackage[margin=24mm,headheight=15pt]{geometry}
\usepackage{amsmath,amssymb,amsthm,mathtools}
\usepackage{booktabs,tabularx,array,longtable,enumitem}
\usepackage{xcolor,tikz}
\usetikzlibrary{arrows.meta,calc,positioning}
\usepackage{microtype}
\usepackage[unicode,colorlinks=true,linkcolor=blue!45!black,citecolor=blue!45!black,urlcolor=blue!45!black]{hyperref}
\usepackage{fancyhdr}
\pagestyle{fancy}
\fancyhf{}
\fancyhead[L]{\small General type surfaces and Lagrangian fibrations}
\fancyhead[R]{\small v25}
\fancyfoot[C]{\thepage}
\renewcommand{\headrulewidth}{0.3pt}
\setlength{\parindent}{1em}
\setlength{\parskip}{0.25em}
\setlength{\emergencystretch}{2em}
\setlist{leftmargin=1.8em,itemsep=0.2em,topsep=0.3em}
\theoremstyle{plain}
\newtheorem{theorem}{定理}[section]
\newtheorem{proposition}[theorem]{命題}
\newtheorem{lemma}[theorem]{補題}
\newtheorem{corollary}[theorem]{系}
\theoremstyle{definition}
\newtheorem{definition}[theorem]{定義}
\newtheorem{remark}[theorem]{注意}
\newtheorem{problem}[theorem]{問題}
\newtheorem{example}[theorem]{例}
\renewcommand{\proofname}{証明}
\renewcommand{\refname}{参考文献}
\renewcommand{\contentsname}{目次}
\newcommand{\PP}{\mathbb P}
\newcommand{\CC}{\mathbb C}
\newcommand{\RR}{\mathbb R}
\newcommand{\ZZ}{\mathbb Z}
\newcommand{\QQ}{\mathbb Q}
\newcommand{\TT}{\mathbb T}
\newcommand{\cX}{\mathcal X}
\newcommand{\eps}{\varepsilon}
\newcommand{\dd}{\,\mathrm d}
\newcommand{\figeight}{(S^1\vee S^1)}
\newcommand{\Bl}{\operatorname{Bl}}
\newcommand{\Cone}{\operatorname{Cone}}
\newcommand{\Sk}{\operatorname{Skel}}
\newcommand{\PD}{\operatorname{PD}}
\newcommand{\Diff}{\operatorname{Diff}}
\newcommand{\id}{\operatorname{id}}

\hypersetup{pdftitle={一般型曲面のLagrangian torus fibration研究ノート第25稿},pdfauthor={}}
\title{一般型曲面のLagrangian torus fibration\\研究ノート 第25稿\\[0.5em]\large 四rayの完全符号分類と二単位枝の整合、\\三高非負係数の端点・退化面}
\author{}
\date{2026年10月28日}
\begin{document}
\maketitle
\noindent
研究再開時には、bundle直下の\texttt{00\_CONTINUATION\_POLICY\_ja.txt}を必ず先に読む。
本ノートは、再利用する補題の量化、今後役立つ計算、避けるべき推論、および具体的な次の問題をまとめる。
公開用論文の証明を履歴として複製せず、現在の研究に必要な形で改訂する。
\tableofcontents

\section{現在の到達点と、残っている問題}

対象は滑らかな$X_d\subset\PP^3$、$d\ge5$、Kähler form $[\omega]=aH$である。
fibrationは連結なfibersを持つproper continuous mapで、底は自由な辺のない有限分岐2複体、正則fiberはLagrangian $T^2$である。
開分岐辺のfiberは$S^1\times\figeight$である。
通常の閉曲面を底とするalmost toric fibrationとは設定が異なり、後者にはLeung--Symingtonの分類が適用される\cite{LeungSymington}。
以下の五次の分母障害では，全平面SNC cores，標準$K_4$ cones，$K_5$ branch graphを保ち，
追加のbranch singularitiesを入れない作用適合atlasを対象とする．

\begin{center}
\begin{tabularx}{\linewidth}{@{}p{.27\linewidth}X@{}}
\toprule
主題 & 現在使える結論と制限\\
\midrule
相対SNC構成 & 成分のgenusを制限しない。正のweighted Kähler表示、各成分$n_a\ge2$、指定marked Morse葉層を仮定。normal holonomyは構成中に選べる。\\
次数分割 & $d=\sum_i e_i$、$k\ge3$の任意分割で、$k$ cores、$d\sum_{i<j}e_i e_j$ nodesを構成。全正則fiberがLagrangian。\\
Kähler slice & 指定vertical affine sliceをweighted sharp simplexとして決定。surgeryに使うのは十分小さいray。\\
整数位相 & 正genusのReeb cyclesと複数三重点を含む整数SDRおよび$H_2$基底。$b_2(B)=p_g(d)-\sum_i p_g(e_i)$。\\
matroidal障害 & 三価リンクの$b_1\le8$。$d\ge7$では全ての次数分割のcoreに局所位相障害。\\
五次・六次 & 全平面分割で大域作用適合atlasを構成。mixed五次は任意分母でも接続不可。六次は全分母一・各bigon二nodesに強制。\\
五次の高分母 & 必要支持は三型・二十六標識、最小次数三以上、三点連結。高分母全体の最大公約数は一。$G_{5,14}$ の両単位枝は node 配分 $\{4,1,0\}$，中央の四高 core の正 matching も固定される。$G_{5,16}$ では五node配分も残る。三型の接着と一般分母剛性は未解決。\\
$K_4$の接着 & 二・四・六例外を全排除。二例外はhubの流れで矛盾し、四例外の全六候補は判別式$-47$の局所障害で矛盾する。一つの全低coreがあれば全分母一。\\
三高・四高の正値性 & 全 pair が非零の整数四ray core の正符号は $2K_2$ または $K_3\sqcup K_1$。正 weights の存在は完全 matching と同値。道型 $P_4$ は全排除。\\
三高非負係数 & 固定 $(s,t)$ から全係数を有限復元。$\delta=\pm1$ は全排除。$\delta=0$ は $x=y=0,s<t,0\le n\le z$ の零軸に限られ、固定 $(s,z)$ で有限復元し，端点は固定 $s$ だけで有限。退化面は固定 $s$，その $x\ge1$ は固定 $b$ でも有限。$h\ge b$ かつ $\gcd(h,b-1)=1$ を排除。一般不存在は未解決。\\
core置換 & meridian模型のhomotopy型・交点格子・境界を計算。閉じたcorner circuitsにregular torus bundlesを加える四pants置換も不可能。open corridors等は残る。\\
解析的課題 & special Lagrangian性、canonical Kähler--Einstein collapse、任意Kähler classへの拡張は未解決。\\
\bottomrule
\end{tabularx}
\end{center}

coreは適切に選んだlogarithmic Weinstein skeletonのsymplectic imageである。
Evans--Mauriの弱いstratificationを用い、Whitney条件や特定のRSTZ skeletonを指定していない\cite{EM,RSTZ,Zhou}。
曲面についての結論を、そのまま$\PP^4$の一般型三次元超曲面に移してはいけない。

\section{任意genusのSNC compactificationを再利用するためのデータ}

\subsection{仮定と、固定するものの順序}

$M$をsmooth connected projective surface、$D=\bigcup_{a=1}^rD_a$をsmooth componentsからなるreduced SNC divisorとする。
三重交点はなく、異なる二成分の複数交点は許す。
$g_a=g(D_a)$、$n_a=\#(D_a\cap\bigcup_{b\ne a}D_b)\ge2$とし、
\[
 [\omega]=2\pi\sum_aw_a[D_a],\qquad w_a>0
\]
を仮定する。
各punctured componentには、境界でradial、内部に$2g_a+n_a-2$個のsimple saddlesだけを持つmarked Morse葉層を指定する。

定理が保つのは、$D_a$上の実際のmarked foliation、crossing product germs、指定したflat disks、および構成途中で選んだcore近傍のprimitive germである。
\emph{normal connectionと正のradiusにおけるcollarのliftは、period correctionの後に選ぶ。}
固定済みの全collar fibrationやholonomyまで保存する追加条件は含まない。
この量化がpositive genusで決定的である。

個々の$D_a$がample、$D_a^2>0$、intersection matrixが可逆、$H^1(M)=0$である必要はない。
正のweighted classがKählerであることから$D$の連結性が従い、正の有理重みへの近似から同じsupportのample divisorが得られるので$M\setminus D$はaffineである。
global logarithmic potentialはweighted curvatureの和と$\dd\dd^c$-lemmaで作る\cite[Lemma (5.11), Remark (5.14)]{DGMS}。
normal crossingのsymplectic regularizationは\cite{TMZ}の枠組みと整合するが、固定複素構造でのKähler flatteningは公開稿のpotentialによる局所議論を使う。

\subsection{normal couplingと全critical leafの保存}

circleの周期は$2\pi$、$\iota_{X_H}\omega=-\dd H$、
$\dd^c\log|z|^2=\dd\theta$とする。
punctured portion $P_a\subset D_a$上で
\[
 \Omega_a=\omega_a+\dd(\mu_a\alpha_a),\qquad
 \lambda_a=(\mu_a-w_a)\alpha_a+\beta_a,\qquad
 \dd\beta_a=\omega_a+w_a\kappa_a
\]
と書く。
$\kappa_a=-\omega_a/w_a$、$\beta_a=0$を選ぶと
\[
 Z\mu_a=\mu_a-w_a,\qquad Z^{\rm horizontal}=0
\]
であり、figure-eight全体と全regular leavesが保存される。
曲線におけるLagrangian liftのcurvature条件にgenus制限がないことは、既に
Evans--MauriのRemark 2.23に明示されている\cite{EM}。
再利用する追加部分は、finite end、指定foliation、全periods、core germの相対的制御である。

面積とnormal degreeの関係は
\[
 \operatorname{area}(D_a)
 =2\pi\left(w_aD_a^2+\sum_{p\in D_a\cap D_b,\ b\ne a}w_b\right),
\]
\[
 \int_{D_a}\kappa_a
 =\frac{2\pi\sum_{p\in D_a\cap D_b}w_b-\operatorname{area}(D_a)}{w_a}
 =-2\pi D_a^2.
\]
従って負のself-intersectionや複数交点も同じ計算で扱える。

\subsection{全periodsを消すholonomy補正}

$G$をdual multigraph、$Q$をintersection matrix、$S=\partial N(D)$とすると、
整数係数で
\[
0\longrightarrow\operatorname{coker}Q
\longrightarrow H_1(S;\ZZ)\longrightarrow H_1(D;\ZZ)\longrightarrow0,
\]
\[
0\longrightarrow\bigoplus_aH_1(D_a;\ZZ)
\longrightarrow H_1(D;\ZZ)\longrightarrow H_1(G;\ZZ)\longrightarrow0.
\]
torsionを含むmeridian relationsは第一列に保持される。

component上のclosed form $\eta_a$をcrossingとflat disksの近傍で零に選び、
\[
 \alpha_{a,s}=\alpha_a+s\pi^*\eta_a
\]
とする。
$\iota_{Y_s}\Omega_{a,s}=-\mu_a\pi^*\eta_a$で定めるMoser flowはzero sectionとnormal first jetを固定し、$\mu_a$を保存する。
すると
\[
 [\Phi_s^*\lambda_{a,s}-\lambda_{a,0}]
 =-s\,w_a[\pi^*\eta_a].
\]
従って各$D_a$の$2g_a$個のperiodsを全て制御できる。
$\Omega_{a,s}$は補助modelであり、引き戻し後は$\Omega_{a,0}$に戻るので、actual Kähler formをこの操作で変更しない。
残る$H^1(G)$のperiodsはcap arms上のbasic formsで消す。

これを省いてbasic Reeb formsだけで補正すると、正genusでは一般に不足する。
\[
 \dim H^1(D)=2\sum_ag_a+e-r+1,\qquad
 \dim H^1(\Gamma)=\sum_ag_a+e-r+1.
\]
holonomyを固定する場合には$\sum_ag_a$次元の不足が残り得る。
「positive genusは障害である」と無条件に結論するのも、「graph periodsだけで十分」とするのも誤りである。

meridian residuesを合わせた後の差はbounded exact form $\dd f$にできる。
\[
 \lambda=\lambda_0+\dd(\chi(\psi_0)f)
\]
というcutoffでは、$Z\psi_0$に$\chi'$の項が現れず、endでの
$O(\rho^{-2})$の正項が$O(\rho^{-1})$の補正を支配する。
このためcore primitive germとnonescaping setを保持できる。

\subsection{finite endpointの式}

transition capでは
\[
 \lambda=(\mu-w_a)\dd\theta_a+
 (I-w_b+\mu A_{ab}(I))\dd\theta,
\]
\[
 A_{ab}=0\ \text{near crossing},\qquad
 A_{ab}=(w_b-I)/w_a\ \text{near }P_a.
\]
ここで
\[
 ZI=\frac{I-w_b+w_aA_{ab}(I)}{1+\mu A'_{ab}(I)}
\]
であり、$ZI=I-w_b$と書いてはいけない。
$I\in[a,b]$、$a>0$、$|ZI|\le C$、$\mu_0\le\varepsilon<w_a/2$なら
\[
 T_a=\log\frac{w_a}{w_a-\mu_0}\le\frac{2\varepsilon}{w_a},
 \qquad |I(t)-I_0|\le\frac{2C\varepsilon}{w_a}.
\]
$\varepsilon<a w_a/(4C)$で別divisorへの先着を排除できる。
ここで$|ZI|\le C$はbufferを付けた拡大intervalで取り、上側bufferからも出ないように$\varepsilon$を縮める。
$C=0$ならこの制約は不要である。以下のendpoint微分は$\mu_0$一定のlink上で用いる。
到達値と単調性は
\[
 I_{\rm end}=\frac{w_aI_0-\mu_0w_b+w_a\mu_0A_{ab}(I_0)}{w_a-\mu_0},
 \quad
 \frac{\partial I_{\rm end}}{\partial I_0}
 =\frac{w_a(1+\mu_0A'_{ab}(I_0))}{w_a-\mu_0}>0.
\]
二到達時刻の最小値を使うのは、二つのactual normal momentsを持つproduct crossing領域だけである。

\subsection{使いやすい試験例と追加条件}

$\PP^1\times\PP^1$のtoric boundaryはvalence二、個々には非ampleであり、linkは円周となる。
二つのsmooth $(1,1)$ curvesなら二交点を持ち、
\[
 Q=\begin{pmatrix}2&2\\2&2\end{pmatrix},\quad
 H_1(S;\ZZ)=\ZZ^2\oplus\ZZ/2,\quad \chi(\Sigma)=2.
\]
これを追加の解析なしに「二つのfocus points」と同定してはいけない。
二つのsmooth plane cubicsが九点で交わる例では
\[
 g_1=g_2=1,\quad b_1(\Gamma)=10,\quad |V_3(\Gamma)|=18,\quad
 \chi(\Sigma)=12.
\]
この例は、両成分がpositive genusを持つ場合のperiod補正を調べる具体的な試験例となる。

$n_a\ge2$は、境界でconstantな実数値Morse functionに内部extremaを持たせないという今回のモデルに必要な仮定であり、一般の葉層に関する不存在定理ではない。
また固定$\omega$で非直交に交わる$D$に、commuting toric normal rotationsを要求することはできない。
定理は同じKähler classの$\omega'$を作る。
元の$\omega$へglobal Moserで戻すときは、$D$自体もisotopyにより移り、元の$D$を集合として保存するとは限らない。

重みを変える問題では、同じ$\omega$のprimitive間のresidue difference $w'-w$が可能であるためには
$\sum_a(w'_a-w_a)[D_a]=0$が必要である。
core primitive germを同じまま保つ追加条件の下では、今回のretractionの議論により各重みも固定される。
cohomology上の自由度とcore相対の自由度を分ける。

\section{次数分割：構成と整数位相の共通書式}

$d=e_1+\cdots+e_k$、$k\ge3$とし、一般のdegree $e_i$ hypersurfaces $S_i=\{F_i=0\}$とdegree $d$ polynomial $G$を用いて
\[
 \prod_iF_i+tG=0
\]
を考える。
$S_2=\sum_{i<j}e_ie_j$、$S_3=\sum_{i<j<\ell}e_ie_je_\ell$と略記する。
各$C_{ij}$は実際の同じmarked curveとして両側で扱う。
\[
 g_{ij}=1+\frac{e_ie_j(e_i+e_j-4)}2,\qquad
 n_{ij}=e_ie_j(d-e_i-e_j)\ge d-2,
\]
\[
 v_{ij}=2g_{ij}+n_{ij}-2=e_ie_j(d-4),\qquad
 a_{ij}=3g_{ij}+n_{ij}-3.
\]
順序付き小解消は
\[
 D_i=\operatorname{Bl}_{d e_i\sum_{j>i}e_j}S_i,\qquad
 N_{\rm node}=dS_2.
\]

\subsection{Picard群とHamiltonian neck}

$i<j$、$Q_{ij}=C_{ij}\cap\{G=0\}$とすると
\[
 N_{C_{ij}/D_i}=\mathcal O_{C_{ij}}(e_j)(-Q_{ij})
 =\mathcal O_{C_{ij}}(e_j-d),\qquad
 N_{C_{ij}/D_j}=\mathcal O_{C_{ij}}(e_i).
\]
triple-point divisor $T_{ij}^{\rm div}$のline bundleは
$\mathcal O_{C_{ij}}(d-e_i-e_j)$なので
\[
 N_{C_{ij}/D_i}\otimes N_{C_{ij}/D_j}
 \otimes\mathcal O_{C_{ij}}(T_{ij}^{\rm div})\simeq\mathcal O_{C_{ij}}.
\]
positive genusではdegreeの和が零だけではPicard群の同定を与えない。
ここでは実際の$F_i,G$がline-bundle同定を供給する。

punctured curveはSteinで$H^2=0$だからholomorphic normal bundlesを自明化できる。
Stein tubular neighborhood上の実際のequationのunitを吸収して$xy=t$を得る。
normal rotationsを平均化した$\Omega$について$\iota_V\Omega$はclosedで、zero sectionへの制限が零である。
tubeがzero sectionへretractすることからexactとなり、Hamiltonian momentが存在する。
genus自体はこの手順を妨げない。

\subsection{weighted Kähler simplex}

\[
 \xi=aH+\sum_i\beta_i[D_i],\quad
 S=\sum_i e_i\beta_i,\quad a_i=a+S-d\beta_i,\quad
 E_\ell=\sum_{i\le\ell}e_i.
\]
\[
 \xi_i=a_iH_i-\sum_{j>i,p}(\beta_j-\beta_i)E_{ij,p}.
\]
sharp条件は
\[
 \beta_1<\cdots<\beta_k,\qquad A=a_k>0,
\]
すなわち
\[
 \delta_\ell=\beta_{\ell+1}-\beta_\ell>0,\qquad
 \sum_{\ell=1}^{k-1}E_\ell\delta_\ell<a.
\]
\[
 \xi_i\cdot C_{ij}=e_ie_j a_j=\xi_j\cdot C_{ij}\quad(i<j),
 \qquad \sum_i\xi_i^2=da^2.
\]
非boundary curveについてBézoutから$\xi_i\cdot\widetilde C\ge A(H_i\cdot C)>0$となる。
実classは$\xi/a$に正規化した後、有理点の凸結合を使う。
symplectic surgeryには$\beta_i=(i-k)\varepsilon$の十分小さいrayだけを使う。

\subsection{底空間とcoreの数値}

\[
 \#C_0=k,\quad \#C_1=(d-4)S_2,\quad
 \#C_2=\sum_{i<j}(3g_{ij}+n_{ij}-3)+S_3,
\]
\[
 \chi(\Sigma_i)=e_i(d^2-4d+6-S_2),\qquad
 \sum_i\chi(\Sigma_i)+dS_2=d(d^2-4d+6).
\]
内部Reeb graphでspanning treeを選ぶと、整数SDRのtargetは
\[
 C_*(\Sk_2\Delta^{k-1};\ZZ)
 \oplus\ZZ^{\,\sum_{i<j}g_{ij}+\sum_{i<j<\ell}(e_ie_je_\ell-1)}[2].
\]
extra generatorsは、non-tree内部辺のcycleと、同じ三成分の異なるtriple pointsのtriangle差である。
したがって
\[
 B_{\boldsymbol e}\simeq
 \bigvee^{p_g(d)-\sum_i p_g(e_i)}S^2,\qquad
 p_g(e)=\frac{(e-1)(e-2)(e-3)}6\quad(e\ge1).
\]
公開稿には符号付きpath chains $D(p)$による$q,s,H$、整数基底、全恒等式の証明がある。

\begin{center}
\begin{tabular}{@{}cclrr@{}}
\toprule
$d$ & 分割 & coreのEuler数 & nodes & $b_2(B)$\\
\midrule
5 & $(1,2,2)$ & $(3,6,6)$ & 40 & 4\\
6 & $(2,2,2)$ & $(12,12,12)$ & 72 & 10\\
6 & $(1,1,4)$ & $(9,9,36)$ & 54 & 9\\
7 & $(1,2,4)$ & $(13,26,52)$ & 98 & 19\\
\bottomrule
\end{tabular}
\end{center}
同じ六次曲面でも、分割を変えると底空間のhomotopy typeを変えられる。
これはcore数だけを変える操作より強い区別である。

\section{matroidal core links：位相分類と作用座標}

\subsection{全次数に対する障害と五次の同時分類}

simple connected rank-three matroid の三価リンクは、二価点を除くと
\[
\begin{array}{c|c|c|c}
|E|&M&\text{link}&b_1\\\hline
4&U_{3,4}&K_4&3\\
5&M(K_4-e)&K_{3,3}&4\\
6&M(K_4)&\text{Petersen}&6\\
7&F_7&\text{Heawood}&8
\end{array}
\]
だけである。非連結 matroids と Euclidean factors も含めた上界は $b_1\le8$。
Fano の体上の realizability は要求しない。nested graph の
$2t_i+s_i=|E|-1$、$t_i+s_i\le3$ が $|E|\le7$ を与える。
既知の $K_{3,3}$ と Petersen の例および nested-set triangulation は
Feichtner--Sturmfels\cite{FeichtnerSturmfels}に基づく。

次数分割の第 $i$ core では
\[
 b_1(\Gamma_i)=1+\frac{e_i(d-e_i)(d-4)}2\ge\binom{d-2}{2}.
\]
$d\ge7$ では、どの分割・marked Reeb graphs を用いても全 core が非 matroidal である。
$d=6$ でも非線形成分は排除され、全平面退化の六つの同時 Petersen 系だけが候補として残る。
六個という数は $K_6$ の one-factorizations の数である。
これらは現在の core を保つ場合の障害であり、core surgery を許す一般の fibration を排除しない。

五次の全 core の同時 matroidal 化は、次のように決まる。
\[
\begin{array}{c|c|c}
\text{分割}&\text{全 core の同時実現}&\text{同時に可能な最大 core 数}\\\hline
(1,1,3)&\text{不可}&2\\
(1,2,2)&\text{不可}&2\\
(1,1,1,2)&\text{可：四つの標識付き抽象系}&4\\
(1,1,1,1,1)&\text{可：一つの抽象系}&5
\end{array}
\]
再利用する構造は、divisor component の内部 Reeb graph $H_a$ が core link の
\emph{連結な誘導部分グラフ}になることである。
\[
 |V(H_a)|=2g_a+n_a-2,\qquad b_1(H_a)=g_a,
 \qquad n_a=3|V(H_a)|-2|E(H_a)|.
\]
simple target の girth が $\ell$ で $g_a>0$ なら、$2g_a+n_a-2\ge\ell$ が必要である。
$(1,1,3)$ の genus-one・三 punctures の block は triangle または多重辺を含み、
三次成分の $K_{3,3}$ を妨げる。両 plane cores を $K_4$ にすると cubic link は triangular prism になる。

$(1,2,2)$ では、四つの共通 triple labels 上の二つの quartet splits を
$\sigma_{AB},\sigma_{AC}$、elliptic four-cycle の bipartition を $\tau$ とすると、
\[
 \Gamma_A=K_4\Longleftrightarrow\sigma_{AB}\ne\sigma_{AC},\qquad
 \Gamma_B=K_{3,3}\Longleftrightarrow\sigma_{AB}=\tau,
 \quad\Gamma_C=K_{3,3}\Longleftrightarrow\sigma_{AC}=\tau.
\]
これが同時実現を妨げる。ただし任意の二 core は同時に実現できる。

$(1,1,1,2)$ では三つの plane--quadric quartets がそれぞれ二つの triple-point pairs を
横断して分ける必要がある。その二択を $\varepsilon_A,\varepsilon_B,\varepsilon_C\in\ZZ/2$ とすると、
quadric link が $K_{3,3}$ となる条件は
\[
 \varepsilon_A+\varepsilon_B+\varepsilon_C=1\pmod2.
\]
偶数なら triangular prism。三点集合ではなく、各 double curve 上の個々の triple points を保持する。
四個という数は finite graph systems の数であり、embedded marked isotopy classes の数ではない。
Morse 実現には good orientation が必要\cite{MichalakReeb}で、公開稿では必要な小 graph に
明示的な高さ関数を与え、actual curves 上へ移している。

\subsection{五直線 Petersen core の作用座標}

五直線 $L_a$、$a\in\ZZ/5$ に正の weights $w_a$ を与え、
$[\omega]=2\pi\sum_aw_a[L_a]$ とする。
各 $L_a$ の quartet を $\{a-1,a+1\}\mid\{a-2,a+2\}$ とすれば Petersen link を得る。
その二 saddle vertices を $A_a,B_a$ とする。

各 puncture $L_a\cap L_b$ に接する外側 Reeb-edge 領域の面積を $2\pi w_b$ にする。
四つの領域の面積の和は全体より $2\pi w_a$ 小さいので、内部 annulus $E_a$ は
\[
 \operatorname{area}(E_a)=2\pi w_a
\]
を持つ。この面積指定は、crossing disks を固定して正の density を調整し、
相対 Moser isotopy で foliation の代表を動かすことにより、指定した marked isotopy class 内で実現できる。
固定した actual critical leaves を点ごとに保って面積を変える主張ではない。

proportional-curvature collar の puncture framing では
\[
 h_{a;b}=\frac{w_b-\operatorname{area}(D_{a;b})/(2\pi)}{w_a}=0.
\]
従って critical ray 上で tangential meridian の action は零、normal action は負となる。
cycle を $\mu_a$、action を $I_{\mu_a}$ と書き、radial moment と区別する。
crossing face $vw$ で $v\in L_a,w\in L_b$ なら
\[
 F=-I_{\mu_a}r_v-I_{\mu_b}r_w.
\]
内部 matching face では
\[
 \beta_a=\mu_{a-1}+\mu_{a+1}
       =-\mu_a-(\mu_{a-2}+\mu_{a+2}),
\]
\[
 F=(-I_{\mu_a}-I_{\beta_a})r_{A_a}+I_{\beta_a}r_{B_a}.
\]
後者の basis change は determinant $-1$。両端で $I_{\beta_a}=0,-I_{\mu_a}$ となり、
全ての branch で $F=-I_{\mu_a}r_v$ に一致する。
十五面の unimodularity、$q=1$ の saturated branch lattice、
Liouville homogeneity と作用比の単調性から、core 近傍全体の標準 Petersen chart を得る。
contact link からの cone construction は Evans--Mauri\cite{EM}に基づく。

重要な論証上の区別は、$\lambda$ が全空間で closed でないことである。
global homology relation に直接 $\int\lambda$ を適用してはいけない。
上の framing relation は、実際の annular torus の $H_1$ で先に確認する。
この torus の $H_1\to H_1(W_5)$ は injective なので、meridian relation をその確認に使える。
その後に、Lagrangian torus 上で closed な $\lambda$ を積分する。

この局所存在を、後述の共有面積・framing・finite-$t$ neck接着と組み合わせると、
六つの全coreを通る大域atlasが得られる。
$M_{0,5}$ filling への置換でもない。五直線側の common-circle determinant は五辺で $0$、十辺で $1$、
標準 $\overline M_{0,5}$ 側では十五辺全てで $1$ であり、boundary $H_1$ の rank も $10$ と $11$ で異なる。

\subsection{固定 period の算術と ambient extension}

固定 annulus が一つの unimodular face となり、両端の branch lattices が saturated なら、
\[
 \frac{\operatorname{area}(E)}{2\pi w}=\frac1{qq'}=\frac1m,
 \quad q,q'\text{ は互いに素な奇数},\quad m\text{ は正の奇数}.
\]
これは同一 torus basis での Farey determinant $\pm1$ と branch の奇数分母条件を合わせたものである。
逆に一つの edge と二つの branch の算術条件には、任意の奇数 $m$ で解がある。
これは geometric periods の同時実現ではない。
unit $p\bmod q$ を二つの units の和にする分解数は
\[
 \prod_{\ell^e\parallel q}\ell^{e-1}(\ell-2),
\]
また互いに素な ordered denominator factorizations は $2^{\omega(m)}$ 個である。

ambient fan lattice は $N=\ZZ^{E(K_4)}/\ZZ\mathbf1$、rank 五である。
正則 torus の $H_1$ は action \emph{cotangent} lattice であり、face ごとに dualization が必要である。
一つの branch block では rank-three integral cotangent extension を作れる。
しかし全 meridians の relation $\sum_a\mu_a=0$ を保つ
\[
 H_1(W_5;\ZZ)\longrightarrow N^*
\]
を同時に要求すると、前項で構成した chart には存在しない。
未定値は五周期方程式 $x_{a-2}+x_{a+2}=1$ により全て $1/2$ に強制される。
rational extension は一意で、その五つの fractional classes は $(\ZZ/2)^4$ を生成する。
それらを含む最小 cotangent overlattice の index は $16$ である。
これは追加した global ambient-linear 条件への障害であり、構成済みの facewise chart を否定しない。
rank 四と五の比較だけでも導けない。

\subsection{四直線 \texorpdfstring{$K_4$}{K4} core と五次の共有ラベル}

四直線補集合では $N_4=\ZZ^4/\ZZ\mathbf1$ とし、四つの labels の有向四周期 $f$ を選ぶ。
$L_i$ の三つの basins に
\[
 \frac{\operatorname{area}(D_{i;j})}{2\pi}
 =w_j+w_i\mathbf1_{\{j=f(i)\}}\qquad(j\ne i)
\]
を割り当てる。全て正で、和は $\sum_jw_j$ である。
相対 Moser により実現でき、critical slope は $h_{ij}=-\epsilon_{ij}$、
$\epsilon_{ij}=\mathbf1_{\{j=f(i)\}}$ となる。
各 crossing face で
\[
 \begin{pmatrix}I_{\mu_i}\\I_{\mu_j}\end{pmatrix}
 =\begin{pmatrix}-1&\epsilon_{ji}\\\epsilon_{ij}&-1\end{pmatrix}
   \begin{pmatrix}t_i\\t_j\end{pmatrix},
 \qquad F=t_ir_i+t_jr_j
\]
とする。四周期には二辺 cycle がないので、六面全ての determinant は $1$。
branch denominator も $1$ となり、一つの $K_4$ core の作用適合 chart を得る。
Evans--Mauri の既存の $K_4$ fibration\cite{EM}に、basin 面積と作用格子の明示式を加えた構成である。
さらに
\[
 \Phi:H_1(W_4;\ZZ)\xrightarrow{\sim}N_4^*,\qquad
 \Phi(\mu_i)=e^*_{f^{-1}(i)}-e_i^*
\]
は整数同型で、上の作用積分と整合する。四周期の incidence map の kernel が対角格子、
image が係数和零の整数格子であることが理由である。
Petersen chart で残る ambient extension の半整数障害は、この $K_4$ chart には生じない。

五次の全平面分割の五 core labels を $\ZZ/5$ とする。第 $i$ core で
\[
 f_i(j)=\frac{i+j}{2}=3(i+j)\pmod5,\qquad j\ne i
\]
と取れば、$j-i$ 上では乗法 $3$ が働くので各 $f_i$ は四周期となる。
しかも $f_i(j)=f_j(i)\ne i,j$ である。
従って共有曲線 $C_{ij}$ の両側で、追加面積を割り当てる basin の triple-point label が一致する。
この共有ラベル系に、後述の canonical weights と $(3,1,1)$ の nodal allocation を組み合わせる。
標識付き対称な四周期系は六つあり、$K_6$ の one-factorizations と対応する。
抽象的な $K_4$ graph system 一個という数とは、数えているデータが異なる。

% Independent notes section. No references to public-paper theorem numbers.
% Existing bibliography keys only: EM, DMT, IKM, VuNgoc.
\section{作用適合 atlas：共有曲線の接着と次に使う整数条件}

\subsection{現在の到達点と適用範囲}

一つの core の作用適合 cone chart は，各正則 face の作用格子を fan の格子に同定し，
branch lattice を飽和させる．大域 atlas では，さらに各 transition と逆写像が
\emph{ambient integer affine map の制限}であることを要求する．
face ごとの unimodularity だけでは，この追加条件は従わない．
$H_1(T^2;\ZZ)$ は作用の\emph{余接線}格子で，fan の $N$ は接線格子である．
focus--focus values は faces の内部に残し，その有限集合を除いた底に atlas を与える．
punctured neighborhoods では実際の作用 monodromy を用いる\cite{VuNgoc}．

現在の比例曲率 SNC compactification の cores と branch graph，標準 coarse cones を保ち，
追加の branch singularities を入れない範囲では，分母を制限せずに大域作用適合 matroidal atlas が得られる
次数分割は\emph{全平面の五次と六次だけ}である．
五次には対称な標識付き四周期 recipes が六つ，六次には Petersen systems が六つある．
どちらも $K_6$ の one-factorizations と対応するが，
この有限個数を非同型な atlases の個数と解釈しない．
七次以上の全分割と六次の混合分割には局所 link 障害があり，
五次の $(1,1,3)$，$(1,2,2)$ は同時 link 条件で，$(1,1,1,2)$ は四 puncture の分母剛性と大域係数で排除される．
core の置換や追加の branch singularities を許す全 fibrations の分類ではない．

\subsection{全平面の五次・六次：共通の weights と node 配分}

小解消の順序を $0<\cdots<d-1$ とし，作用積分の単位を
$I_\gamma=(2\pi)^{-1}\int_\gamma\lambda$ とする．
補助 vertical class $[\Omega]=aH+\sum_i\beta_i[D_i]$ に対し
\begin{equation*}
 b_i=\frac{\beta_i}{2\pi},\qquad
 c=\frac{a}{2\pi(d-1)},\qquad
 R_i=(d-1)c+\sum_jb_j-db_i,\qquad
 w_{i,j}=c+b_j-b_i
\end{equation*}
と置く．row sum は $\sum_{j\ne i}w_{i,j}=R_i$ である．
$i<j$ の共有曲線では $\delta_{ij}=b_j-b_i>0$，
$R_i-R_j=d\delta_{ij}$ であり，$i$ 側の $d$ 回の blow-up 後の実面積は $2\pi R_j$ になる．
全 weights が正となる条件は $\max b-\min b<c$ で，surgery 用の十分小さい vertical ray では満たせる．
これは選んだ residue recipe の条件であり，全 Kähler cone の条件ではない．

五次では labels を $\mathbb F_5$ とし，
$f_i(j)=(i+j)/2=3(i+j)$，
$\epsilon_{ij;k}=\mathbf1_{\{k=f_i(j)\}}$ と取る．
同じ値が $f_j(i)$ からも得られる．puncture $k$ の blow-down basin と node 数は
\[
 B^0_{i;jk}=w_{i,k}+\epsilon_{ij;k}w_{i,j},\qquad
 n_{ij;k}=1+2\epsilon_{ij;k},
\]
\[
 B^0_{i;jk}-n_{ij;k}\delta_{ij}
 =w_{j,k}+\epsilon_{ij;k}w_{j,i}>0.
\]
各 curve の三 basins への配分は $(3,1,1)$ であり，
各 triple の三つの pairs のうち selected pair は一つだけなので，各 triangle に五 nodes が入る．

六次では，両側で同じ quartet を用い，四つの external basins に各一 node，
内部 annulus に二 nodes を置く．共有面積の等式は
\[
 w_{i,k}-\delta_{ij}=w_{j,k}\quad(k\ne i,j),\qquad
 w_{i,j}-2\delta_{ij}=w_{j,i}.
\]
これらの面積条件と row sums を固定すると，五次・六次とも上の $w_{i,j}$ が唯一の解である．
固定された blow-up points は，marked isotopy によって指定領域の相異なる regular leaves に置き，
その小近傍を固定した density と relative Moser で実面積を合わせる．

\begin{center}
\begin{tabular}{c|c|c|c|c}
次数&曲線上の nodes&triangles&bigons&全 nodes\\\hline
5&$(3,1,1)$&$10\times5$&なし&50\\
6&$(1,1,1,1;2)$&$20\times3$&$15\times2$&90
\end{tabular}
\end{center}
全 branch の正規化された affine length は
$w_{i,j}+w_{j,i}=2c=a/(\pi(d-1))$ である．六次では各 pair の二本の branch が同じ長さを持つ．

\paragraph{Cycle framing と star transitions.}
共通 neck の normal cycle を $\alpha$ とすれば反対側は $\alpha'=-\alpha$ である．
一つの blow-up は $x=zu$ により puncture lift に $+\alpha$ を加え，
一つの triple end $uvw=t$ は $-\alpha$ を加える．従って一 basin の実際の torus 内で
\[
 \beta'_k=\beta_k+(n_{ij;k}-1)\alpha.
\]
五次の face transition は，branch/transverse coordinates において
\[
 (s,t)\longmapsto(-s+\theta_kt,t),\qquad
 \theta_k=\mathbf1_{\{f_i(k)=j\}}+\mathbf1_{\{f_j(k)=i\}},
 \quad\sum_k\theta_k=2.
\]
この和が $K_4$ の ray relation を保ち，三枚の面は一つの $\operatorname{GL}(3,\ZZ)$ map で貼れる．
六次の external cycles は両端で同じである．quartet の一方 $K$ を囲む cycle を $\beta$ とすると，
$A$ branch で $\beta'=\beta$，$B$ branch で $\beta'=\beta+2\alpha$ となる．
内部 face に二 nodes があるため，二式を一つの monodromy-free trivialization の式として扱わない．
各 branch の outward matching cycle を使えば，Petersen star の transition は
\[
 r_v\mapsto-r'_v,\quad r_x\mapsto r'_x,\quad r_y\mapsto r'_y,
 \quad r_z\mapsto r'_z-2r'_v.
\]
$r_x+r_y+r_z=r_v$ が保存され，star の飽和 rank-three 基底で determinant は $-1$ である．

\subsection{有限 \texorpdfstring{$t$}{t} の neck 接着を再利用する条件}

中央曲線の面積一致だけを，有限 $t\ne0$ の作用切片の一致へ読み替えない．
使うのは，actual neck の全空間の tube 上の一つの primitive と，固定 moment 上の輸送である．
tube は punctured curve へ retract して $H^2=0$ となる．circle 平均した primitive は，
中央の固定集合によって $\lambda_{\rm neck}(V)=\mu$ と正規化される．
各 toric end の第二 reference action は
$R_\ell(\mu)=\nu_\ell|_{\rm cut}+c_\ell^{\rm ref}$ で，定数は両側と全 $t$ に共通である．
各 reference circle から regular circle までの reduced annulus $A$ では
\[
 J_{\beta_\ell}-R_\ell(\mu)=\frac1{2\pi}\int_A\omega_{t,\mu}.
\]
framing の変更は整数倍の $\mu$ を加える．critical value は regular annuli の exhaustion で取る．
中央の $\mu=0$ では framing 項が消え，共通 reference actions と strict double curve 上の
実際の critical-region areas が定数項を決める．従って共有面積の一致と cycle shear による勾配の一致から，
両成分の critical action germs は同じ affine functions になる．
それを有限 $t$ へ保つ段階には，次の固定切断面での輸送を使う．

固定した $\mu=\delta>0$ と $\mu=-\delta<0$ の open buffers を先に選び，
horizontal transport $\mathcal P_t$ を行う．同じ primitive に関して
\[
 \mathcal P_t^*(\lambda_{\rm neck}|_{X_t})
 -\lambda_{\rm neck}|_{D_i}=\dd G_t
\]
なので，closed transported cycles の作用は切片まで保存される．
$\mu=0,t=0$ を通る輸送や，$\delta\to0$ に一様な輸送は使わない．
boundary reference actions と固定 winding に Stokes を適用すると，
全 reduced area は $2\pi\sum_\ell\pm(R_\ell(\mu)+m_\ell\mu)$，$m_\ell\in\ZZ$ という affine function になる．

再利用する十分条件は，両端の open germs が同じ marked Reeb tree を持ち，
共通 framing と reference actions で
$J^{\rm crit}_{\beta_\ell}=c_\ell+p_\ell\mu$，$p_\ell\in\ZZ$，が一致することである．
この式が指定する各 critical region の target area $a_h(\mu)$ には
\[
 \sum_h a_h(\mu)=\int_C\omega_{t,\mu},\qquad
 a_h(\mu)>\int_{D_h\cap K}\omega_{t,\mu}
\]
を要求する．$K$ には保存する caps と collars 全体を含める．
端で一致する marked interpolation に，parameter に相対的な density--Moser を施すと，
実際の reduced forms の下でこれらの面積を持つ family が得られる．
circle lift は Lagrangian であり，critical actions は全 interval で上の affine functions のままである．
同じ marked graph だけでは measured data の一致は保証されない\cite{DMT,IKM}．
追加の compact parameters を使うときは，初期 interpolation と region identifications も滑らかに与える．
Liouville cone を用いる局所構成の基礎は \cite{EM} にあるが，
固定された二つの end systems を無条件に同一視する主張ではない．

\subsection{法円作用の tropical divisor と core ごとの総和}

三価 star の relation $u_1+u_2+u_3=\kappa v$ と common normal action $h$ に対し
\[
 \mathfrak d(h)=\sum_{a=1}^3\dd h_a(u_a)-\kappa\dd h(v)
\]
を使う．これは単位 facet weights に対する tropical principal divisor の係数であり，
公式と lift 非依存性は Allermann--Rau\cite[Definition 3.4, Remark 3.5]{AllermannRau}にある．
各 core の別々の normal actions を一つの大域関数と同一視しない．
本構成で使うのは，ambient affine transitions で係数が保存され，
同じ neck の逆向き normal circles に対応する二端では相殺するという条件である．

primitive ray parameter $s$ において
\[
 I_{\alpha_v}=-q_vs,\qquad I_{\mu_w}=-p_{v\to w}s,
 \qquad q_v>0,\quad p_{v\to w}\in\ZZ
\]
と書く．三 arm cycles の numerators とその和が $q_v$ と互いに素であるため，$q_v$ は奇数である．
現在の positive sectors における crossing determinant と basin area は
\[
 q_vq_w-p_{v\to w}p_{w\to v}=1,
 \qquad B_{v\to w}=w_w-w_v\frac{p_{v\to w}}{q_v}>0.
\]
一 core の全 vertices の分岐係数は
\[
 \sum_v\mathfrak d_v
 =\sum_v\bigl(\kappa_v-\deg_{\rm int}(v)\bigr)q_v
   -\sum_{v\to w\,\mathrm{crossing}}p_{v\to w}.
\]
各 vertex の incoming numerators を，総和を取った後に outgoing の添字へ入れ替えた式である．
三 puncture の component は basins が全曲線を尽くすため，自己交点を $e$ とすると
$\sum p=-e q$ となる．

\subsection{共有する四 puncture 曲線の分母剛性}

再利用する条件は，actual neck の両端で同じ marked $2+2$ Reeb tree を使い，
二本の branch の primitive action lattices を保ち，追加の branch singularities を入れないことである．
両端で比例曲率 collars と Liouville cone の作用正規化を用いる．
正 normal circles は $\alpha^+=-\alpha^-$ なので，共有 branch の分母は両端で同じ $q_A,q_B$ となる．

自己交点を $e^\pm$，各端の二つずつの puncture numerators の和を $P_A^\pm,P_B^\pm$ とする．
内部 annulus の両 rays の正の格子 index を $m^\pm$ とすれば，actual framing
$\beta_A+\beta_B=-e\alpha$ と面積の正値性から
\[
 e^\pm q_Aq_B+q_BP_A^\pm+q_AP_B^\pm=m^\pm,
 \qquad
 \frac{\operatorname{area}(E^\pm)}{2\pi w^\pm}
 =\frac{m^\pm}{q_Aq_B}>0.
\]
actual torus cycles の整数 clutching は
\[
 \beta_\ell^+=\beta_\ell^-+k_\ell\alpha^-,\qquad k_\ell\in\ZZ,
 \qquad
 \frac{p_\ell^-+p_\ell^+}{q_{s(\ell)}}=-k_\ell.
\]
従って両側の annulus equations の和は
\[
 \frac{m^-+m^+}{q_Aq_B}
 =e^-+e^+-\sum_{\ell=1}^4k_\ell\in\ZZ.
\]
これが再利用すべき整除条件
\[
 \boxed{q_Aq_B\mid m^-+m^+}
\]
である．両 cones が unimodular なら $m^-=m^+=1$ なので，奇数性から $q_A=q_B=1$ となる．
証明は作用の微分だけを比較し，加法定数，node 数の非負性，全 node budget に依存しない．
puncture の degree-one markings と normal circle を保つことが必要である．

ordered resolution の meridian nodes では $k_\ell=n_\ell-1$ である．
単純 node の整数 shear は \cite[\S4.3, equation (4.10)]{SymingtonFour}の convention と対応する．
従って外側四 basins の node 総数は
\[
 N_{\rm ext}=4+e^-+e^+-\frac{m^-+m^+}{q_Aq_B}.
\]
unimodular な仮想接着に必要な数は次の通りである．mixed の大域 atlas は後述の係数で排除される．
\begin{center}
\begin{tabular}{c|c|c|c|c}
共有曲線&$(e^-,e^+)$&全 nodes&external 合計&internal annulus\\\hline
mixed plane--quadric&$(4,2)$&10&8&2\\
全平面 sextic&$(1,1)$&6&4&2
\end{tabular}
\end{center}
個々の external basin への配分まで一意とは述べない．
六次では全十五共有曲線に適用できるため，三十 branches の全分母は一に強制される．

\paragraph{高分母を許すための変更.}
cone index を緩めても $m^-+m^+\ge q_Aq_B$ が必要である．
例えば $(q_A,q_B)=(1,3)$ に二つの index-one annuli を合わせることはできないが，
indices $(1,2)$ はこの整除条件だけは満たす．後者は geometric chart の存在を意味しない．
Reeb tree の変更，複数 cones への展開，orbifold lattice，追加 branch singularities を許すなら，
その変更を明示して actual integral cycles の比較から条件を立て直す．

\subsection{Mixed quintic の任意分母に対する障害}

$(1,1,1,2)$ の各 core は個別には chart を持つ．限定した正の面積 recipes では，
quadric の $K_{3,3}$ が $40/64$，二直線と conic の $K_4$ が $4/36$ で実現する．
しかし三本の conic necks に四 puncture 剛性を適用すると，六本の plane--quadric branches の分母は一となる．
三本の plane--plane branches の分母は任意のままでよい．

mixed plane の line contribution は $\kappa=-1,e=1$ によって
$-q-\sum p=0$ と消える．conic の二端では $q_A=q_B=1$，
$P_A+P_B=1-4=-3$ なので，plane core の総和は $-4-(-3)=-1$ となる．
quadric の各 component は $\kappa=0,e=2$，$P'_A+P'_B=-1$ なので，寄与は $-2-(-1)=-1$ である．
全四 cores の総和は
\[
 3(-1)+3(-1)=-6
\]
となり，九本の branches の対向端での相殺と矛盾する．
従って現在の cores，branch graph，標準 coarse cones を保つ限り，任意分母で大域 atlas は存在しない．

孤立した annulus の整数解は接着の解ではない．例えば $q_A=1,q_B=3$ で
$e^-=4,(P_A^-,P_B^-)=(-1,-8)$ と
$e^+=2,(P_A^+,P_B^+)=(0,-5)$ は各々の index-one 式を満たすが，
外側 node 総数は $28/3$ となり，整数 clutching に反する．
この例は，局所の係数の符号を変えられることだけでは大域障害を回避できないことを示す．

\subsection{全平面五次へ進むための局所奇数分母模型}

四直線の $K_4$ core には三 puncture の曲線だけがあり，四 puncture 剛性は適用できない．
正の weights $w_i$ と整数 $q_i>0,p_{ij}$ が
\[
 \sum_{j\ne i}p_{ij}=-q_i,\quad q_iq_j-p_{ij}p_{ji}=1,
 \quad B_{ij}=w_j-w_i p_{ij}/q_i>0
\]
を満たせば，これらは局所 chart の十分条件となる．
basin areas を相対 Moser で実現し，crossing face の作用行列
\[
 \begin{pmatrix}-q_i&-p_{ji}\\-p_{ij}&-q_j\end{pmatrix}
\]
を使う．各 collar と cap で正の sector を保つことが必要であり，行列式だけを確認して終えない．

任意の $h\ge1$ に対して $q=2h+1$ とし，
\[
 (q_A,q_B,q_C,q_D)=(q,1,1,1),\qquad
 (w_A,w_B,w_C,w_D)=\left(q,h,1,1+\frac1{4h}\right),
\]
\[
 (p_{ij})=
 \begin{pmatrix}
 0&-2h&-2&1\\
 -1&0&0&0\\
 -h&h-1&0&0\\
 2h&0&-q&0
 \end{pmatrix}
\]
とすれば，六つの determinants は一で全 basin areas が正となる．
各 component の normalized area は $W=3h+3+1/(4h)$ であり，共通の正の倍率は自由である．
各 branch の係数は $(-3h,h,2h+2,-2)$ である．
これは一つの core の幾何学的な存在族であり，五つの cores を接着した atlas ではない．

\paragraph{高分母が一本だけ孤立する配置は除外する.}
一般に一つの $K_4$ core の分母が $(q,1,1,1)$，$q>1$ で basin areas が全て正なら，
高分母 vertex の outgoing numerators は順序を除き
\[
 (-(q-1),-(b+1),b),\qquad b\mid q-1,\quad b+1\mid q-1.
\]
$q-1=t b(b+1)$ と書くと，その分岐係数は
$-t(b(b+1)+1)<0$ である．三 numerators が全て負という候補は，
残る三つの $q=1$ vertices に正の directed cycle を強制し，basin positivity に反する．
従って全平面五次の $K_5$ branch graph で，$q>1$ の部分 graph は一本の edge だけからなる
連結成分を持てない．一般の leaf を排除したわけではない．
隣接する複数の高分母 branches については，次節の局所表示，素数 support，
第三 corner 条件を合わせて，端ごとの係数相殺と actual basin areas を調べる．

\section{五次の高分母：局所分類と大域整合}
\label{nt19:denominators}

この節には，今後の高分母の構成・排除に再利用する条件を残す．
高分母を持つ全平面五次の大域作用 atlas の存在・不存在は，なお一般には未解決である．
現在の必要候補は五頂点全てを使う三型・二十六標識であり，
最小次数三以上，三点連結で，高辺は八本以上である．
Book と $G_{5,15}$ は完全に排除済みであり，二価頂点は存在しない．
同じ core に二本の単位枝が接すれば全分母一となる．
非空高支持では単位枝は高々二辺の matching であり，高分母全体の最大公約数は一である．
高分母が存在するなら最大値は百三以上で，少なくとも二種類の高分母が現れる．
この有界排除と，分母上限を用いない graph 障害とは量化が異なる．

三高分母の正符号 graph は完全 matching であることが証明されたため，
固定低データの二次復元と canonical Schur 線形系を全ての正整数三高模型に使える．
三高係数の同時非負性は，$b=-\beta\ge2$，$d=-D_0\ge1$，
$\alpha\ge k\ge b+1$ の非特異領域に限られる．
低 row の二整数 $(s,t)$ だけを固定した全係数・全分母の有限復元を保持する．
本版では $\delta=\pm1$ を全排除し，比例領域 $\delta=0$ を
$x=y=0,z\ge1,s<t,0\le n\le z$ の零軸へ帰着した．
この零軸は固定 $(s,z)$，その端点は固定 $s$，退化面 $a=0$ も固定 $s$ で完全有限復元できる．
退化面の $x\ge1$ には互いに素な場合の障害と固定 $b$ の完全有限復元もあるが，
各固定整数自体への一様上界と三高同時非負性の一般不存在は未解決である．

三高・四高に共通して，全 pair が非零なら正 weights の存在は正符号 graph が
$2K_2$ であることと同値となる．道型 $P_4$ は正値恒等式により全域で排除された．
三高 core 間の単位枝は node 配分 $\{5,0,0\}$ または $\{4,1,0\}$ に限られ，
隣接分母の少なくとも一つは三または五である．
$G_{5,14}$ の両単位枝では五node配分も排除され，ともに $\{4,1,0\}$ に限られる．
唯一の四高 core の正 matching も，二つの単位枝の端点同士を結ぶものに固定される．
この構造と外側 pair の復元を使って，残る自由尺度を大域条件へ接続する．

\subsection{正値性と二高分母の表示}

局所四頂点に $q_i>0$，$p_{ij}\in\ZZ$ を置く．
\[
 \sum_{j\ne i}p_{ij}=-q_i,\qquad
 p_{ij}p_{ji}=q_iq_j-1,\qquad
 D_i=-q_i-\sum_{j\ne i}p_{ji}.
\]
正の basin weights の存在は，$p_{ij}>0$ の directed edges からなる全 simple cycles で
\[
 \prod_{i\to j}p_{ij}<\prod_iq_i
\]
となることと同値である．二閉路は determinant 条件で自動的に通るので，
最大八個の有向三角形と六個の有向四角形を検査すればよい．
これは既知の差分不等式・max-times の判定\cite[\S2]{SergeevSchneiderButkovic}である．
全 gains に有理数 $\theta>1$ を十分一に近く掛け，終点 $j$ の simple paths の gain と一の最大を
$w_j$ と取れば，正の有理 weights の証明書を構成できる．

二高分母 $(a,b,1,1)$，$a,b>1$ では，high vertices および low vertices の交換を許すと，
次の表示が全 numerator matrices を尽くす．
\[
 r>0,\quad r\mid a-1,\quad r+1\mid b-1,\quad
 u\ne0,\quad u\mid a-1,\quad
 x=\frac{a-1}{r},\quad y=\frac{b-1}{r+1},
\]
\[
 A=a+x+u,\quad A\mid ab-1,\quad
 B=\frac{ab-1}{A},\quad v=B-b+y\ne0,\quad v\mid b-1.
\]
このとき $A,B>0$ であり，
\[
 P=\begin{pmatrix}
 0&-A&x&u\\
 -B&0&-y&v\\
 r&-r-1&0&0\\
 (a-1)/u&(b-1)/v&d&0
 \end{pmatrix},\qquad
 d=-1-\frac{a-1}{u}-\frac{b-1}{v}.
\]
正の weights が存在するための必要十分条件は
\[
 u<0<v
\]
と分かっている．算術稿の \texttt{ls11:two-high-signs} は $a=b$ を仮定しない．
従って正の模型を探索するとき，$u>0$ の候補を残す必要はない．
各 high row のもう一つの high への entry は負であり，二つの low への entries は一正一負となる．
固定した $(a,b)$ の局所幾何は，この有限な divisor 計算と局所 chart の十分条件で判定できる．
分母対 $(a,b)$ の全体が有限個という主張ではない．

奇数 $b>1$ に対し，$(3,b,1,1)$ の正の局所 chart が存在することと
$b\in\{7,11\}$ は同値である．この結論にも $b$ の上限はない．
非正の numerators だけなら全 $q_i=1$ なので，通常の arithmetical structures の
非正非対角行列\cite[\S\S1--3]{CorralesValencia}へ高分母全体を帰着しない．

\paragraph{一般二高分母の負係数と，再利用する整数変数．}
正の標準形で $u=-s<0$，$v=z>0$，$S=(a-1)/s$，$T=(b-1)/z$ と置くと，
$d=S-T-1<0$ であり，$\min(D_A,D_B)<0$ も一般に成立する．
証明は算術稿の \texttt{td12:low-entry} と \texttt{td12:one-negative} にある．
二高分母が同時に非負の $D$ を持つ候補を再探索する必要はない．
再利用する変数は
\[
 h=S-r,\quad k=T-r-1,\quad
 \alpha=x-s=sh/r,\quad \beta=y-z=zk/(r+1)
\]
であり，これらは全て非零で同符号となる．さらに
\[
 M=rs+1,\quad N=(r+1)z+1,\quad
 \delta=\alpha-\beta,\quad L=M-N+r\delta
\]
と置けば
\[
 M\beta-N\alpha=1,\qquad \alpha L=1+a\delta,
 \qquad D_A=h-L,\quad D_B=L+\delta-k
\]
となる．特に $\gcd(\alpha,\beta)=1$ である．
負枝 $\alpha,\beta<0$，$\delta\geq0$ では，整数性が $S=T$ を強制する．
この枝では $d=-1$ かつ $D_B<0$ となる．

正の局所 weights は
\[
 (w_A,w_B,w_C,w_D)
 =\left(a,b,x+\frac1{2r},z+\frac1{2T}\right)
\]
と明示できる．正の entries に対応する四つの basin areas は
$1/(2r),1/2,1/(2T),1/2$ となるため，整数データから面積正値性の証明書を直接作れる．
この定理は分母の上限，奇数性，大域 clutching，node budget を用いない．
ただし一般の三高分母で全ての高係数 $D$ が非負となることを禁止できるかは未解決である．
三高分母の numerators の matching 性とは別の問題である．

\paragraph{再試行を避けるための符号反例．}
一意の最大分母の係数が負とは限らない．二高分母だけでも
\[
 q=(19,17,1,1),\qquad
 P=\begin{pmatrix}
 0&-23&6&-2\\-14&0&-4&1\\3&-4&0&0\\-9&16&-8&0
 \end{pmatrix},\qquad
 w=\left(19,17,\frac{37}{6},\frac{33}{32}\right)
\]
は正の局所模型であり，$D=(1,-6,5,0)$ となる．
従って閉路に沿う $D$ の向きから分母の単調性を無条件には結論しない．
三高分母の既存例 $(13,7,7,1)$ でも $D=(4,-2,-4,2)$ となる．
また，異なる二高分母の係数和も常に負ではない．
具体的には
\[
 q=(11,1,1,9),\qquad
 P=\begin{pmatrix}
 0&-2&5&-14\\
 -5&0&0&4\\
 2&-1&0&-2\\
 -7&2&-4&0
 \end{pmatrix},\qquad
 w=(33/16,1,1,33/8)
\]
は正の局所 chart であり，$D=(-1,0,-2,3)$ を与える．従って
\[
 \sum_{q_i>1}D_i=2>0,\qquad
 \sum_i\frac{D_i}{q_i}=-\frac{58}{33}<0.
\]
一方，等二高分母 $q=7$ の標準模型では $D=(0,-2,2,0)$ であり，
$\sum_iD_i/q_i=12/7>0$ となる．したがって，単純な高分母部分和や
この重み付き和の一様な符号を，大域の potential として仮定しない．

\paragraph{等三高分母の負和．}
分母 $(q,q,q,1)$，$q\geq3$，では低分母頂点の係数は $D_0>0$ であり，
高分母三頂点の係数和は負となる．この結果には正値性を用いない．
正局所模型では一般の matching 定理からも，低分母 row が二負一正であることが分かる．
その row を $(-a,-b,a+b-1)$ と書けば
\[
 D_0=-1+(q-1)\left(\frac1a+\frac1b-\frac1{a+b-1}\right)>0
\]
となる．正値性を仮定しない元の等三高分母定理は，別途 row/pair の整数条件から証明されている．
等二高分母の解決済みの分類と，それに使った整数変数を次にまとめる．

\input{notes_uniform_two_high.tex}

\subsection{単色の高分母成分と二値分母の排除}

高分母 support を $H=\{ij:q_{ij}>1\}\subset K_5$ とする．
辺を持つ一つの連結成分 $C$ の全高分母が同じ $q>1$ なら，
各 core の分母 pattern は $(q^k,1^{4-k})$ となる．
$k=1,2,3$ では高分母側の $D$ の和は負であり，$k=4$ では全四頂点の和なので零である．
しかし $C$ 上で足すと各 edge の対向端が相殺し，総和は零でなければならない．
従って全 vertices の次数が四となり，$C=K_5$ となるが，これは素数 support の
$K_4$ 禁止に反する．

従って，高分母 support の各非空連結成分には少なくとも二種類の高分母が必要である．
特に，ある $q>1$ に対して全 branch denominators が $\{1,q\}$ に属すれば，全て一となる．
算術稿の \texttt{uh11:monochromatic-component} と
\texttt{uh11:two-value-rigidity} に対応する．
この障害は分母上限，node 非負性，全 node budget を用いない．
「単色の部分 graph」が別の高分母 edge に接する場合には，連結成分という仮定が失われるので，
その部分 graph だけを同じ議論で排除しない．

\subsection{素数 support と共通因子の制約}

五 cores の labels を $V=\{0,1,2,3,4\}$ とする．
共有 branch の分母を $q_{ij}=q_{ji}$，core $i$ の outgoing numerator を $p^i_{jk}$ と書く．
\[
 \sum_{k\ne i,j}p^i_{jk}=-q_{ij},\qquad
 p^i_{jk}p^i_{kj}=q_{ij}q_{ik}-1,\qquad
 q_{ij}\mid p^i_{jk}+p^j_{ik}.
\]
奇素数 $\ell$ ごとに $G_\ell=\{ij:\ell\mid q_{ij}\}\subset K_5$ とする．
その三角形 $ijk$ 上では，$\mathbf F_\ell$ において
\[
 p^i_{jk}=p^j_{ki}=p^k_{ij}=\varepsilon,\qquad
 p^i_{kj}=p^j_{ik}=p^k_{ji}=-\varepsilon,
 \qquad\varepsilon=\pm1.
\]
三つの determinants と三つの clutchings を一周すると $x^2=1$ を得る．
この式と row sums から $K_4$ と book $B_3$ が禁止され，$G_3$ は matching となる．

現在の密な support に対する一般結論は
\[
 |E(G_\ell)|\geq7
 \quad\Longrightarrow\quad
 \ell=5,\quad G_\ell\cong K_{2,2,1}
\]
である．従って exact 七辺 support はどの奇素数でも存在せず，
\[
 |E(G_3)|\leq2,\qquad |E(G_5)|\leq8,\qquad
 |E(G_\ell)|\leq6\quad(\ell\ne3,5).
\]
五の場合も edge 数七は不可能である．七辺 graph の残存二型
$K_1\vee P_4$ と $K_{2,3}$ に一本を加えた graph は，
全 cores の式を同時に使うと八本目の辺を強制する．
後者では最後の determinant の残差が $-5$ になることが，素数を五へ限定する核心である．

四 cores に対する以前の gcd 条件に加え，任意に選んだ七本の異なる branches に対して
\[
 \gcd\{q_{ij}:i,j\in S,\ i<j\}=1\quad(|S|=4),
 \qquad
 \gcd(q_{e_1},\ldots,q_{e_7})\mid5
\]
となる．特に二十五で割れる分母は高々六本である．
$G_5=K_{2,2,1}$ の場合，その八辺の分母の gcd は正確に五であり，
少なくとも二辺では五進付値が一となる．
この指数の障害は $\ZZ/25\ZZ$ 上で証明する．
その環では $r^2=0$ から $r=0$ としてはならず，nilpotent ideal を保持して
最後の determinant 残差 $-5\ne0$ を取り出す．

固定された標識付き $K_{2,2,1}$ に対し，modulo five の reduced system は符号 $s=\pm1$ の二解だけである．
parts を $\{a,b\},\{c,d\},\{o\}$ とすると
\[
 q_{ab}=-s,\quad q_{cd}=s,\quad q_{ij}=0\ (ij\ne ab,cd),
 \qquad
 P^o=s\begin{pmatrix}
 0&-2&1&1\\-2&0&1&1\\-1&-1&0&2\\-1&-1&2&0
 \end{pmatrix}
\]
となり，他の四 cores も一意に決まる．これらは対向 $D$ の相殺も満たすため，
同じ mod~$5$ 条件だけでこの型をさらに排除できない．
正の奇整数への lifting，全整数 clutching，面積，node 配分，actual neck の接着は別に確認する必要がある．

三に関する局所計算も再利用できる．一つの core で二つの分母が三で割れると仮定すると，
それらの rows は $\mathbf F_3$ 内で一定符号となり，残る二 rows の最初の二 entries は打ち消す．
最後の pair は $p_{23}=-q_2$，$p_{32}=-q_3$ を要求し，determinant の $-1$ と矛盾する．
実際の atlas で全ての高分母が三で割れるなら，$H=G_3$ が matching となり，
孤立高分母 edge の障害から $H$ は空である．
素数 support の合同定理には weights や node 非負性を使わないが，
この最後の幾何的結論には正の basin areas と対向 $D$ の相殺を使う．

\subsection{低分母 puncture と三角形成分の一般障害}

同じ高分母 $q$ の共有 branch の一つの marked puncture が，両側で低分母一へ向かうとする．
その outgoing entries $p,p'$ は $q-1$ の非零 signed divisors であり，
$q\mid p+p'$ から
\[
 p'=-p,\qquad\{p,p'\}=\{1,q-1\},\qquad
 \{p,p'\}=\{-1,-(q-1)\}
\]
の三型に限られる．各型で incoming entries の和も $p+p'$ に等しい．
meridian node 表示での個数は順に $1,0,2$ である．
これは marked punctures の対応を固定して使う必要がある．

$H$ の一つの連結成分が三角形なら，各 core は二高分母を持つ．
二つの外側 low punctures の outgoing sums は，符号が揃う場合に $q,-q$，
逆に対応する場合に零，零となるため，総和は常に零である．
従って三角形に向かう high puncture の outgoing numerators の和は $-2q$ となり，
対向 $D$ の相殺も使うと high--high numerators の差は一つの circulation $\delta$ に一致する．
それらの絶対値の平均を $T_i>0$ と書けば
\[
 q_i=\frac{T_i+T_{i+1}}2,\qquad
 (T_{i-1}+T_i)(T_i+T_{i+1})=4T_i^2-\delta^2+4.
\]
最大値と最小値の比較から $\delta^2=4$ と全 $T_i$ の一致を得る．
従って三辺分母が等しくなり，単色成分の障害に反する．
局所の追加計算では $q=7,13$ の high--high pair $(-q-1,-q+1)$ だけが残り，
その高分母係数は $(0,-2)$ または $(0,-1)$ で，やはり相殺できない．

この三角形成分の排除は分母の上限によらず，node 非負性や全 node 数五も仮定しない．
三角形に別の高分母 edge が接続する場合は，この議論の範囲に含まれない．
一般の path 成分には，次に述べる別の局所符号の議論を用いる．

\input{notes_fixed_low_completion_v17.tex}
\input{notes_support_v25.tex}

\subsection{第三 corner の有限条件と一辺で整合する無限族}

二つの cores $i,j$ を共有辺 $ij$ で比較し，puncture $k$ に対して
\[
 A=q_{ik},\quad B=q_{jk},\quad r=p^i_{kj},\quad s=p^j_{ki}
\]
が既知とする．第三 core $k$ の二つの numerators は
$e=Au-r$，$f=Bv-s$ と書かれ，$u=1-n_{ik;j}$，$v=1-n_{jk;i}$ である．
各共有直線の全 node 数五と非負性を使うと，必要条件は
\[
 (Au-r)(Bv-s)=AB-1,
 \qquad u,v\in\{-4,-3,-2,-1,0,1\}
\]
となる．この三十六候補の検査は分母の大きさによらず使える．
三 punctures 全てで通っても，第三 core の全データや actual weights の同時実現はまだ必要である．

再利用できる局所無限族として，$n\geq1$，$Q_n=n^2+n+1$，$R_n=n^2-n+1$ に対し
\[
 q=(Q_n,1,1,R_n),\qquad
 P_n=\begin{pmatrix}
 0&-(n+1)&1&-(n^2+1)\\
 -n&0&0&n-1\\
 n(n+1)&-(2n+1)&0&-n(n-1)\\
 -n^2&n&-1&0
 \end{pmatrix},
\]
\[
 w=(Q_n-1/2,1,1,n-1+1/(2n)),\qquad
 D=(-Q_n,2n+1,-1,R_n)
\]
を持つ正局所模型 $L_n$ がある．$L_n$ の $Q_n$ 端と $L_{n+1}$ の $R_{n+1}=Q_n$ 端では
$D$ が相殺し，outgoing および incoming numerators の和は，適切な共通順序で
$(0,0,-2Q_n)$ となる．node 配分は $(1,1,3)$，star shears は $(0,0,2)$ である．
しかし $n\geq2$ では residue matching が一意で，第三 corner は
\[
 (R_nu+n^2)\bigl(Q_{n+1}v+n^2+2n+2\bigr)
   =n^2(n+1)^2+2
\]
を要求し，六値の $u,v$ の全場合が一般の多項式不等式で排除される．
この族は局所模型と単独辺の整数整合の存在例であり，actual weights や切片まで接続した大域例ではない．

\subsection{既存の有界排除の量化と，残る研究の入口}

$q_{ij}\le101$ の全範囲では全 $q_{ij}=1$ が証明済みである．
奇数性から，高分母が存在すれば最大値は百三以上となる．
この結論は分母が全て百一以下という仮定を必要とし，現在の三候補の無上界排除を意味しない．

完全性に必要な手法を残す．二高分母以上の core では $q_0,q_3>1$ と並べ替え，
$p_{01},p_{02},p_{13}$ をそれぞれ $q_0q_1-1,q_0q_2-1,q_1q_3-1$ の signed divisors から選ぶ．
残りは row sums と pair products で決まるため，numerators に独立の探索上限は不要である．
一高分母では完全な高 ray の標準形を保持する．端点の fingerprint
\[
 \left(q,D,\operatorname{sort}(p_1,p_2,p_3\bmod q)\right)
\]
について，対向端での $D$ と residues の符号反転を要求し，相手のない候補を反復削除する．
これは必要条件による削除なので，大域接着の十分条件とはしない．
元の完全 catalog と証明書は
\nolinkurl{verify_quintic_denominators_v10.py} と
\nolinkurl{quintic_denominators_verification_v10.json} に保持し，版番号を付け替えない．

次の対象は引き続き $G_{5,14},G_{5,16},G_{5,17}$ である．
単位枝の小分母・node 表と，三高・四高に共通する完全 matching を使い，
分母の上限を増やす前に全 core の符号と node の位置を整合させる．
Book と $G_{5,15}$ は全排除済みであり，再探索しない．その過程で得た負枝の二直線分類と
小さい $L$ の正枝の全尺度表示は，残る支持の接着へ再利用する．
三高分母の同時非負性では非特異な canonical Schur 系に集中できるが，
一般の局所模型の分類には正整数解を持つ特異系を保持する．
必要な整数候補が得られれば全 core の weights，整数 framings，実作用切片まで検討する．
密な mod 五模型の整数 lifting では，五進付値一の辺が少なくとも二本ある条件を保持する．

$G_{5,5},G_{5,7},G_{5,9},G_{5,10},G_{5,11},G_{5,12},G_{5,13},G_{5,15}$ は全排除済みであり，再探索しない．
その証明から，共通低 ray の $C/D$ 型，matching が強制する符号，
未知分母を切らない点・半直線の処理を取り出して次の候補へ使う．
等二高分母の負和，素数 support の合同障害，または部分的に整合する整数模型を，
任意分母の全体に対する不存在・存在へ無条件に移さない．


\section{core と meridian nodes を取り込む局所置換}

\subsection{meridian blow-ups の再利用できる全データ}

$m\ge4$ 本の一般位置直線の $i$ 番目の成分上で、相異なる smooth points を $r_i$ 個 blow up し、
直線の strict transforms だけを除いた補集合を $U_{\boldsymbol r}$ とする。
$r=\sum r_i$、$s=\#\{r_i>0\}<m$、$k=m-1-s$ と書く。
infinitely near points や任意の vanishing cycles はこの模型に含めない。
Hattori の meridian 標識と整合する模型\cite{Hattori,PapadimaSuciu}に disk を付けると
\[
 U_{\boldsymbol r}\simeq(T^k)^{(2)}\vee\bigvee^tS^2,
 \qquad t=\binom{m-1}{2}-\binom{k}{2}+r-s.
\]
\[
 \pi_1=\ZZ^k,\quad b_2=\binom{m-1}{2}+r-s,\quad
 \chi=\binom{m-2}{2}+r.
\]
compact core の intersection lattice、signature、boundary は
\[
 Q_{\mathrm{abs}}=0^{\binom{m-1}{2}}\oplus\bigoplus_{r_i>0}(-A_{r_i-1}),
 \qquad \sigma=-(r-s),
\]
\[
 Q_{\mathrm{plumbing}}=J_m-\operatorname{diag}(r_i),\qquad
 H_1(Y)=\ZZ^{\binom{m-1}{2}+k}\oplus\bigoplus_{r_i>0}\ZZ/r_i.
\]
reduced JSJ graph は $K_m$。境界包含の kernel は抽象群として
$\ZZ^{\binom{m-1}{2}}\oplus\bigoplus\ZZ/r_i$ であり、具体的な peripheral subgroup の保存は別に必要である。
untouched line があるため ample divisor の support を作れ、$U_{\boldsymbol r}$ は affine となる。
canonical bundle は自明で、対応する Stein boundary では
\[
 c_1(\xi)=0,\qquad d_3(\xi)=\frac{-2\chi+3(r-s)}4
\]
である\cite{Gompf}。

六次の一 core と一つの指定 meridian node の模型は $m=5,r=s=1$ で、
\[
 U\simeq(T^3)^{(2)}\vee3S^2,\qquad
 \pi_2(U)=\ZZ[\ZZ^3]^4,
\]
\[
 H_1(Y)=\ZZ^9,\quad Q_{\mathrm{abs}}=0,\quad
 \chi=4,\quad d_3=-2,\quad \text{JSJ graph}=K_5.
\]
同じ homotopy 型だけでは、この boundary を復元できない。

一般の canonical pants budget $P=(m-3)^2$ を目標にするなら、必要 node 数は
$r=\binom{m-3}{2}$。さらに $\pi_1=\ZZ^3$ を要求すれば $s=m-4$ となり、
\[
 U\simeq(T^3)^{(2)}\vee(P-1)S^2,\quad
 \pi_2(U)=\ZZ[\ZZ^3]^P,\quad
 \sigma=-\binom{m-4}{2}.
\]
full-face tree assembly の intersection form は零なので、$m\ge6$ ではこれだけでも置換できない。
元の次数は $d=m+1$ である。

\subsection{閉じた corner circuits の充填まで含む no-go}

標準四次元 pants の互いに離した full Seifert facets を対で貼る。
その結果の\emph{閉じた} corner circuits に、compact surfaces 上の任意の regular $T^2$-bundles を加えることまで許す。
disk filling、非自明な monodromy、非 disk base も含むが、open corner chains へは付けず、circle collapse や singular fibers は加えない。
external boundary は連結で negative-base Seifert pieces を持つとする。
Golla--Martelli\cite[Proposition 4.3, Remark 4.9]{GollaMartelli}の一般の構成全体とは区別する。

pants 数を $p$、facet pairs の数を $e$、その graph を $\Gamma$ とし、
\[
 \mathcal C(Y)=\sum_{\text{JSJ pieces}}-\chi_{\rm orb}(\text{base})
\]
とする。この class では
\[
 \chi(X)=p,\qquad\mathcal C(Y)=4p-2e.
\]
tree component には closed corner circuit がないので、regular fillers により他 component と接続できない。
従って連結 $X$ は、
\[
 \Gamma\text{ tree},\quad\mathcal C=2\chi+2,
 \quad\text{reduced JSJ graph は planar}
\]
または
\[
 \text{全 components が cycle を持つ},\quad\mathcal C\le2\chi
\]
のいずれかとなる。対象 $U$ は $\chi=4,\mathcal C=10$、JSJ graph $K_5$ なので不可能である。
この証明は完成後の基本群を仮定せず、disk relations を加えても回避できない。

多直線の boundary は $\mathcal C=m(m-3)$ である。
$p=(m-3)^2$ の場合、facet pairs は $e=3(m-3)(m-4)/2$ に強制され、
$\Gamma$ が連結なら $g=b_1(\Gamma)=\binom{m-4}{2}$。
disk-only completion では $\Gamma$ は連結である。
完成後の $\pi_1=\ZZ^3$ に対して、$g>0$ なら disk 数 $q\ge\max\{1,g-3\}$。
vertex の meridian $\ZZ^3$ が完成後も同型に残る\emph{追加条件}を置けば $q\ge g$。
この survival 条件を省いて $q\ge g$ と主張してはいけない。

次の候補には open torus corridors、partial-facet gluing、circle collapse など、上の定義外の操作が必要である。
一般の relative pants surgery の不存在は証明していない。

\subsection{具体的な affine chart と Kirby 比較の入力}

一つの一般位置模型を
\[
 U\cong\CC^2_{t,z}\setminus
 \{(1-tz)(1-(2t+1)z)(1-(3t+4)z)=0\}
\]
と取れる。$E^\circ=\{z=0\}\cong\mathbb A^1$。
$(t,z)\mapsto(tz,z)$ は、三 affine lines の補集合への affine modification である\cite{KalimanZaidenberg}。
全 algebraic units は $E^\circ$ 上で定数なので、$U$ は very affine ではない。

deleted curve 上の $u=t+3z$ は finite degree-six projection である。
三 node values は $-4,-7/2,-18/5$、六 simple branch values は
\[
 \pm2\sqrt3,\quad-\tfrac12\pm\sqrt6,\quad-\tfrac{10}{3},\quad\tfrac23.
\]
全て相異なる。reference fiber を $u=4$ とし、distinguished paths に沿って
三 quadratic equations の roots を運べば attaching braids を決められる。
Sugawara\cite[Proposition 3.1, Theorems 3.7 and 4.1]{Sugawara}により smooth handles の数は
$(h_0,h_1,h_2)=(1,6,9)$ となる。Weinstein attaching data の同定は未完了である。

大円上の roots は原点近くの三点 cluster と $u/3$ 近くの三点 cluster に分かれる。
その full twists を $T_-,T_+$、六本全体を $T_6$ とすると
\[
 \beta_\infty=T_6T_-^{-2}T_+^{-1},\qquad
 \text{exponent sum}=12.
\]
これは roots の大半径 isotopy による等式で、pairwise linking だけから推論していない。
以前の非 proper ruling における四 strand braid $\Delta_4^2$ とは、projection と fiber chart が異なる。
六 branch values と三 curve nodes を、九つの interior focus--focus fibers と解釈してはいけない。

次はこの equation と paths から\emph{framed} Kirby diagram を作り、
open torus corridor を含む pants 側の diagram と、peripheral kernel、五つの Seifert slopes、
contact collar、actual periods を比較する。smooth handle 数の一致だけで surgery を認定しない。


\section{相対gluingで再採用してはいけない近道}

\subsection{実際のmarked curveと、同じHamiltonian moment}

両側で同じactual marked surfaceを用いる。
abstract labelled Reeb graphだけではpure mapping classを決められない。
crossing charts、embedded marked foliations、regular flat patchesを先に選び、
局所許容量を確保してからblow-up size、normal collar、$|t|$を順に小さくする。

\emph{同じactual $\Omega$と同じHamiltonian $\mu$}の上で$h_\mu$を補間する場合、
$\{\mu,h_\mu\}=0$でLagrangian性が得られる。
reduced diffeomorphismがarea-preservingである必要はない。
独立に固定された二つのintegrable systemsをsymplectically同定する問題では、
measured Reeb graphやfocus--focus semi-global invariantsが必要である\cite{DMT,IKM,VuNgoc}。
compact nodal fiber全体とそのbufferを飽和truncationに含め、actual transportする。

\subsection{解消後のopen component全体はexactとは限らない}

明示的平面退化
\[
 M_s=x_0+s x_1+s^2x_2+s^3x_3,\qquad
 \prod_{i=1}^dM_i+t\prod_{r=d+1}^{2d}M_r=0
\]
では
\[
 p_{ijk}=[-ijk:ij+ik+jk:-(i+j+k):1].
\]
最初の成分内の補助直線のstrict transform
\[
 Q_r=H_1-\sum_{j=2}^dE_{1j,r}
\]
はboundaryを避けるcompact curveで、$\int_{Q_r}\omega=a>0$である。
したがって解消後の$D_1^\circ$全体にLiouville primitiveは存在しない。
exact coreとnodal patchesを分ける必要がある。

\subsection{成分上のgermを全smoothingへ延ばすこと}

標準複素平面のHamiltonian
$H(z)=-\frac13\operatorname{Im}(z^3)$のtime-one mapは
$g(z)=z+\bar z^{\,2}+O(|z|^3)$で、原点とfirst derivativeを固定する。
仮に$F=(X,Y)$が$XY=xy$を保つ$C^2$ germで
$F(x,0)=(g(x),0)$なら
\[
 \partial_yY(x,0)=\frac{x}{g(x)}
 =1-\frac{\bar x^{\,2}}x+O(|x|^2).
\]
一次項$-re^{-3\sqrt{-1}\theta}$は実線形写像の一次項ではないので、
$\partial_yY$は原点で微分可能でない。
normal first jetの一致だけで全$xy=t$ germへの延長を結論できない。
actual annular transportがこの問題を避ける。

\section{次に着手する具体的問題}

算術側の現在の入口は \S\ref{ns25:support-agenda} にまとめた．
次の焦点は $G_{5,14}$ の二つの $\{4,1,0\}$ と固定された中央 matching の大域整合である．
四高 $P_4$ と $G_{5,14}$ の五node sector は解決済みなので再探索しない．
$G_{5,16}$ では五node sector が残り，$G_{5,17}$ には単位枝の議論を直接使えない．
三高係数では，比例零軸の端点の一般整数根問題，内部 $0<n<z$，
退化面の $h\ge b$ における非互いに素な領域，一般 $|\delta|\ge2$ を区別する．
退化面で $x\ge1$ なら必ず $h\ge b$ であるが，$x=0,h<b$ にはこの互いに素な場合の障害を適用せず，
既存の固定 $s$ 復元を保持する．
新しい完全有限手順は定理に基づく固定入力ごとの判定に使い，入力上限を増やすだけの探索は優先しない．
反例が現れれば全原式と正 weights を確認し，一般不存在を固定目標とせず構造記述へ転換する．

密な $G_5=K_{2,2,1}$ の二つの合同模型も，整数 lifting の出発点として残る．
少なくとも二本の五進付値一を保ち，正整数性，clutching，node 配分を同時に課す．
分母上限の数値を増やすだけの探索は優先しない．
候補が残れば，五つの cores の actual weights と basin areas，integer framings を合わせ，
最後に有限の平滑化 parameter の作用切片と ambient star transitions を確認する．
合同模型や必要条件の存在を，global atlas の存在と同一視しない．

幾何的な方向として、構成した五次・六次 atlas の変形と不変量を調べる。
六つの標識付き recipes または graph systems が、標識を忘れたときにどのように同値になるかを調べる。
integer affine holonomy、translation data、focus invariants を別々に記録する。
底とfibrationの同値性を調べる不変量と、個々のtorusの埋込みを区別する不変量を分ける。
後者には displacement energy germ を候補として用いる。
$t$、vertical polarization、foliation の選択を変えたときに何が不変かを調べる価値がある。
focus punctures を持つ atlas に適用できる tropical homology や affine sheaf の定式化を先に確認し、
compact smooth tropical manifold の定理を無条件に使わない。
Evans--Urz\'uaのrational homology ballsの埋込み\cite{EvansUrzua}と，
Brendel--Hauber--Schmitzのsmall toriおよびdisplacement energy germ\cite{BHS}は，
局所模型からambient symplectic manifold内のtoriを比較するための資料となる．
適切なsymplectic formに対する存在と，本研究で指定したKähler formでの実現は別に確認する．
個々のtorusの非同値性だけで，底全体のatlasやfibrationの非同値性を結論しない．
Cartwright の ordered weak tropical complex では最初の component の intersection matrix が負定値なので、
今回の action atlas の構造と同一視できない。

局所置換の方向として、open torus corridor を含む具体的な四 pants 置換を調べる。
三 hyperbola equations と finite projection から framed Kirby diagram を作り、pants 側と比較する。
$\pi_2=\ZZ[\ZZ^3]^4$、零交点形式、$K_5$ JSJ boundary と peripheral kernel を同時に追う。
closed-circuit regular completions は任意 monodromy を含めて排除済みである。
その候補を繰り返さず、変更する open corner attachment を先に明示する。
証明は独立した公開稿「直線配置の meridian blow-up と四次元 pants 置換の障害」にまとめた。

一般 SNC 相対定理を別の actual degeneration へ適用する方向も残る。
normal line bundles の同型、component cores、同時 corner transport を検査する。
matroidal 性を緩める場合、Aksnes の arroid fans\cite{Aksnes}は
transverse very affine plane arrangements の比較対象になるが、Reeb 展開や actual periods の同定は追加問題である。
三次元では Stein curve tube が高次元 strata に替わるため、曲面の証明をそのまま使わない。

この順位は暫定的である。面白さとインパクトを最大化するよう、その時点の状況と先行研究を踏まえて
研究計画を柔軟に変更する。研究ノートと公開用論文は必要に応じて大幅に改訂し、
ノートには将来の研究に役立つ内容だけを残す。
公開用論文はトータルの価値を最大化するよう、必要に応じてさらに分割または再編する。
本版では公開稿を三篇に分けた。第一篇は相対SNC構成・大域作用atlas・底の位相、第二篇は五次の局所整数理論と大域接着障害、第三篇は固定境界の局所置換を扱う。各篇は独立して読める導入と必要な定義を持つ。

\section{一次文献の射程と検算}

Mikhalkinのpants decompositionはglobalなstratified fibrationを与えるが、
各最大cellに少なくとも一つのLagrangian fiberという量化を、全正則fiberへ拡張して読まない\cite{Mikhalkin}。
FBPのivy-like fibrationは大きな部分領域の構成であり、全component coreのfillingは追加の問題である\cite{FBP}。
CG K3のtransport定理はKulikov/almost toricの仮定を持ち、改訂Lemma 2.13の極限先は必要なcylindersを潰した商である\cite{CGK3}。
構成・幾何稿のIntroductionでは、これらとfinite SNC構成・大域作用atlasの関係を詳しく引用する。算術稿ではmax algebra、principal minors、arithmetical structures、friezesとの相違を一次文献に基づいて整理する。
別稿ではarrangement topology・plumbing・pants assembly・affine modification・handle constructionの先行研究を独立に位置づける。
法円作用の分岐係数はAllermann--Rauの既知のtropical divisor公式であり，新しい内容はSNCでの計算と対向端比較である\cite{AllermannRau}。
正のweightsの閉路判定は既知のmax-times理論に位置づける\cite{SergeevSchneiderButkovic}．
非対角成分の符号が自由であるため，通常のarithmetical structuresの有限性\cite{CorralesValencia}を分母の上限へ流用しない．

今後追跡する発表として，Mak--Matessi--Ruddat--Zharkov による compact Calabi--Yau
hypersurfaces の大域 fibration の講演概要がある\cite{RuddatAnnouncement2026}．
Liouville flow を用いる混合局所模型の帰納的構成が述べられており，
一般型への相対接着との比較に役立つ可能性がある．
加えて Zharkov の2026年4月15日の講演概要は，
phase/ober-tropical hypersurfaces と soft-cell models を pants の観点から比較する
\cite{ZharkovSoftCells2026}．ページ掲載日は4月8日であり，講演日と区別する．
今回確認したこれらの一次資料は講演概要である．
完成論文の証明を確認したことや，正式 preprint が存在しないことは意味しない．
正確な仮定と相対性の範囲を引き続き調べ，幾何稿の Introduction と構成の位置付けを改訂する．

\subsection{torus fibers の不変量を調べる際の文献上の入口}
\label{nfi24:section}

幾何的な次の課題として，既に構成した五次・六次 atlas の同値関係と，
その torus fibers の不変量を比較する方向を残す．
両者は同じ問題ではなく，一つの torus の非同値性から fibration 全体の非同値性を直ちに結論しない．

Venugopalan--Woodward\cite[Definitions~1.1--1.2, Theorem~1.3]{VenugopalanWoodward2026}の
tropical disk-count formula を参照する場合は，次の仮定を先に確認する．
対象は compact almost toric 四次元多様体とその torus fiber であり，
elementary polyhedral decomposition を持ち，potential に寄与する全 disks の tropical graphs が
generic interior points にのみ collisions を持つことを要求する．
本研究の開分岐 fiber $S^1\times(S^1\vee S^1)$ と core fibers を備えた底複体は，
この定理の設定へ自動的には入らない．局所 torus の比較領域を選ぶなら，
disk の compactness と領域外への逸出を含めて新たに比較を証明する必要がある．
非単調な場合の potential は Maurer--Cartan space 上の関数を同型まで扱うため，
任意の表示の Laurent polynomial を正準不変量として比較しない．

当面の再利用可能な問いは，指定した core・neck に相対的な変形が，
既知の affine holonomy，translation data，focus--focus invariants と
fiber の Floer 不変量にどう反映されるかである．
文献の仮定を本構成へ橋渡しできる見通しが得られれば，
支持 graph の全排除より幾何方向を優先してよい．


\subsection{現在の問題に直結する検算への入口}

\begin{itemize}
\item \nolinkurl{verify_four_ray_matching_v25.py}：原四行・五 pair の有理式から第六 pair の正値恒等式を検査し，全六十四符号 graph を分類する．
\item \nolinkurl{verify_g514_five_node_v25.py}：七局所型，三十六全正因子，恒等零線形式の整除障害を確認する．
\item \nolinkurl{verify_g514_central_matching_v25.py}：中央交差 matching の四正規形と全 node 零配置，元 clutching と行和の矛盾を確認する．
\item \nolinkurl{verify_high_d_axis_endpoints_v25.py}：比例零軸の端点の全係数証明と，固定 $s$ の全有限復元を行う．
\item \nolinkurl{verify_degenerate_divisor_v25.py}：退化面の元残余の積表示，正値性，固定 $b$ の有限復元を確認する．
\item \nolinkurl{verify_unit_three_high_triangles_v24.py} と \nolinkurl{verify_outer_core_transfer_v23.py}：残る単位枝の全三角表と共鳴を保持する．
\item \nolinkurl{verify_support_agenda_v23.py}：全 $1024$ graph，三型・二十六標識を確認する現行支持証明書である．
\end{itemize}

完全復元器の固定データを混同しない．
\nolinkurl{research/exact_three_high_completion_v17.py} は固定低 row/column，
\nolinkurl{verify_high_d_low_row_finiteness_v22.py} \texttt{--fixed s t} は固定低 row の二整数，
\nolinkurl{verify_high_d_axis_degenerate_v24.py} \texttt{--fixed-s s} は退化面の固定 $s$，
\nolinkurl{verify_high_d_proportional_v24.py} \texttt{--axis s z} は比例零軸の固定 $(s,z)$ を扱う．
新しい \nolinkurl{verify_high_d_axis_endpoints_v25.py} の \texttt{--s-cap N} は比例零軸の両端点を，
\nolinkurl{verify_degenerate_divisor_v25.py} の \texttt{--fixed-b b} は退化面 $x\ge1$ を扱う．
それぞれ残りの係数・分母には任意の探索上限を置かない．
旧 \nolinkurl{verify_high_d_asymmetry_v23.py} \texttt{--axis b z} の対象は本版で排除済みであり，
現在の残余軸を調べる復元器ではない．
Book と $G_{5,15}$ の専用検算は証明依存の再現性資料として保持する．

旧検算の source・結果・監査を内容不変で保持し，新しい五主検算を追加する．
現行の支持 graph 証明書は引き続き
\nolinkurl{support_agenda_verification_v23.json} であり，候補は三型・二十六標識のままである．
旧四型・五十六標識は $G_{5,15}$ 排除前の中間段階である．
検算と追加監査の一覧，実行範囲，source と証明の hash は
\nolinkurl{verification/README.txt} と各結果に記録する．
内部の独立点検を外部査読や proof assistant による形式化と表現しない．

全 $q\le101$ の排除は有界完全検算である．
無上界の支持障害，固定 $(s,t)$ の全係数復元，特定例，有限回帰とは量化を区別する．
異なる証明書の重複した候補数を加算しない．
一般命題の証明は公開稿に置き，コードは整数性・恒等式・有限型分解を検査する．
整数係数を全て記録した非負多項式の符号証明と，有界な変数探索を混同しない．
有限検算は Moser，analytic gluing，Lagrangian fibration の存在証明の代用ではない．

\begin{thebibliography}{99}

\bibitem{CGblowup}
P. Chakravarthy and Y. Groman,
\emph{Almost toric fibrations on symplectic blow ups},
arXiv:2510.00994v1 (2025).
Theorem~1.1, Definition~2.2, Lemma~3.7.
\url{https://arxiv.org/abs/2510.00994v1}.

\bibitem{CGK3}
P. Chakravarthy and Y. Groman,
\emph{Almost toric fibrations on K3 surfaces via degenerations},
Advances in Mathematics \textbf{502} (2026), Part~B, 111138.
Definition~1.2, Theorems~1.6, 1.8, Corollary~1.9;
Lemma~2.13 is cited in the revised arXiv version of 28 May 2026.
\url{https://doi.org/10.1016/j.aim.2026.111138};
\url{https://arxiv.org/abs/2502.04304v3}.

\bibitem{DMT}
J.-P. Dufour, P. Molino and A. Toulet,
\emph{Classification des systèmes intégrables en dimension 2 et invariants des modèles de Fomenko},
Comptes Rendus de l'Académie des Sciences, Paris, Série~I, Mathématique
\textbf{318} (1994), no.~10, 949--952.
An English translation is available at
\url{https://webpages.csus.edu/mjv/DufourMolinoToulet.pdf}.

\bibitem{Evans}
J. Evans,
\emph{Lectures on Lagrangian Torus Fibrations},
Cambridge University Press, Cambridge, 2023.
\url{https://arxiv.org/abs/2110.08643v4}.

\bibitem{EM}
J. D. Evans and M. Mauri,
\emph{Constructing local models for Lagrangian torus fibrations},
Annales Henri Lebesgue \textbf{4} (2021), 537--570.
Lemmas~2.16, 2.22, 3.1; Remark~2.19; \S3.1.
\url{https://doi.org/10.5802/ahl.80};
\url{https://arxiv.org/abs/1905.09229v3}.

\bibitem{FBP}
J. Fernández de Bobadilla and T. Pełka,
\emph{Fibrations by Lagrangian tori for maximal Calabi--Yau degenerations and beyond},
Communications on Pure and Applied Mathematics (2026), e70060,
published online 8 July 2026,
\url{https://doi.org/10.1002/cpa.70060}.
Theorem~6.14, Proposition~5.6, Example~5.8 and Remark~5.9 are cited from the final preprint,
arXiv:2312.13248v4, 9 July 2026.
\url{https://arxiv.org/abs/2312.13248v4}.

\bibitem{GollaMartelli}
M. Golla and B. Martelli,
\emph{Pair of pants decomposition of 4-manifolds},
Algebraic \& Geometric Topology \textbf{17} (2017), no.~3, 1407--1444.
\url{https://doi.org/10.2140/agt.2017.17.1407};
\url{https://arxiv.org/abs/1503.05839}.

\bibitem{IKM}
A. Izosimov, B. Khesin and M. Mousavi,
\emph{Coadjoint orbits of symplectic diffeomorphisms of surfaces and ideal hydrodynamics},
Annales de l'Institut Fourier \textbf{66} (2016), no.~6, 2385--2433.
Theorems~3.11 and 3.17 (measured Reeb graphs and weighted Reeb graphs).
\url{https://doi.org/10.5802/aif.3066};
\url{https://arxiv.org/abs/1504.05629}.

\bibitem{LeungSymington}
N. C. Leung and M. Symington,
\emph{Almost toric symplectic four-manifolds},
Journal of Symplectic Geometry \textbf{8} (2010), no.~2, 143--187.
Theorem~2.6 and Table~1.
\url{https://doi.org/10.4310/JSG.2010.v8.n2.a2};
\url{https://arxiv.org/abs/math/0312165}.

\bibitem{Mikhalkin}
G. Mikhalkin,
\emph{Decomposition into pairs-of-pants for complex algebraic hypersurfaces},
Topology \textbf{43} (2004), no.~5, 1035--1065.
Theorems~1 and~1$'$; \S4.2.
\url{https://doi.org/10.1016/j.top.2003.11.006};
\url{https://arxiv.org/abs/math/0205011v3}.

\bibitem{RSTZ}
H. Ruddat, N. Sibilla, D. Treumann and E. Zaslow,
\emph{Skeleta of affine hypersurfaces},
Geometry \& Topology \textbf{18} (2014), no.~3, 1343--1395.
\url{https://doi.org/10.2140/gt.2014.18.1343};
\url{https://arxiv.org/abs/1211.5263}.

\bibitem{VuNgoc}
S. Vũ Ngọc,
\emph{On semi-global invariants for focus--focus singularities},
Topology \textbf{42} (2003), no.~2, 365--380.
Theorem~2.1 and \S3.
\url{https://doi.org/10.1016/S0040-9383(01)00026-X};
\url{https://arxiv.org/abs/math/0208245}.

\bibitem{Zhang}
Y. Zhang,
\emph{Pair-of-Pants Decompositions of 4-Manifolds Diffeomorphic to General Type Hypersurfaces},
Journal of Mathematical Study \textbf{58} (2025), no.~2, 211--252.
Theorem~1.1; the quintic example in \S1.
\url{https://doi.org/10.4208/jms.v58n2.25.06};
\url{https://arxiv.org/abs/2102.10037}.

\bibitem{Zhou}
P. Zhou,
\emph{Lagrangian skeleta of hypersurfaces in $(\mathbb C^*)^n$},
Selecta Mathematica (N.S.) \textbf{26} (2020), no.~2, article~26, 33~pp.
\url{https://arxiv.org/abs/1803.00320}.

% Bibliography for the SNC theorem and subsequent applications.

\bibitem{DGMS}
P. Deligne, P. Griffiths, J. Morgan and D. Sullivan,
\emph{Real homotopy theory of K\"ahler manifolds},
Inventiones Mathematicae \textbf{29} (1975), no.~3, 245--274.
Lemma~(5.11), pp.~266--267, and Remark~(5.14).
\url{https://doi.org/10.1007/BF01389853};
\url{https://www.math.stonybrook.edu/~dennis/publications/PDF/DS-pub-0024.pdf}.

\bibitem{TMZ}
M. Farajzadeh Tehrani, M. McLean and A. Zinger,
\emph{Normal crossings singularities for symplectic topology},
Advances in Mathematics \textbf{339} (2018), 672--748.
Definitions~2.8, 2.9, 2.11, 2.12 and Theorem~2.13.
\url{https://doi.org/10.1016/j.aim.2018.09.035};
\url{https://arxiv.org/abs/1410.0609v3}.

\bibitem{CohenSuciu}
D. C. Cohen and A. I. Suciu,
\emph{The boundary manifold of a complex line arrangement},
Geometry \& Topology Monographs \textbf{13} (2008), 105--146.
\S2.5, Proposition~2.6, Theorem~3.2.
\url{https://doi.org/10.2140/gtm.2008.13.105};
\url{https://arxiv.org/abs/math/0607274}.

\bibitem{Gompf}
R. E. Gompf,
\emph{Handlebody construction of Stein surfaces},
Annals of Mathematics \textbf{148} (1998), no.~2, 619--693.
Definition~4.2 and Theorem~4.5; the invariant denoted $\theta$ there equals $4d_3$ here.
\url{https://doi.org/10.2307/121005};
\url{https://arxiv.org/abs/math/9803019}.

\bibitem{KalimanZaidenberg}
Sh. Kaliman and M. Zaidenberg,
\emph{Affine modifications and affine hypersurfaces with a very transitive automorphism group},
Transformation Groups \textbf{4} (1999), no.~1, 53--95.
\S\S1--3; in the arXiv version, Definition~1.1 and Proposition~3.1.
\url{https://arxiv.org/abs/math/9801076v2}.

\bibitem{Neumann}
W. D. Neumann,
\emph{A calculus for plumbing applied to the topology of complex surface singularities
and degenerating complex curves},
Transactions of the American Mathematical Society \textbf{268} (1981), no.~2, 299--344.
\S5, especially the reduced graph criteria on pp.~321--322;
Theorems~4.1 and~4.2 for plumbing normal forms.
\url{https://doi.org/10.1090/S0002-9947-1981-0632532-8};
\url{https://webhomes.maths.ed.ac.uk/~v1ranick/papers/neumann008.pdf}.

\bibitem{FeichtnerSturmfels}
Eva Maria Feichtner and Bernd Sturmfels,
\emph{Matroid polytopes, nested sets and Bergman fans},
Portugaliae Mathematica (N.S.) \textbf{62} (2005), no.~4, 437--468.
\href{https://arxiv.org/abs/math/0411260}{arXiv:math/0411260}.

\bibitem{ArdilaKlivans}
Federico Ardila and Caroline J. Klivans,
\emph{The Bergman complex of a matroid and phylogenetic trees},
Journal of Combinatorial Theory, Series B \textbf{96} (2006), no.~1, 38--49.
\href{https://doi.org/10.1016/j.jctb.2005.06.004}{doi:10.1016/j.jctb.2005.06.004}.
\href{https://arxiv.org/abs/math/0311370}{arXiv:math/0311370}.

\bibitem{Shaw}
Kristin M. Shaw,
\emph{A tropical intersection product in matroidal fans},
SIAM Journal on Discrete Mathematics \textbf{27} (2013), no.~1, 459--491.
\href{https://doi.org/10.1137/110850141}{doi:10.1137/110850141}.
\href{https://arxiv.org/abs/1010.3967v2}{arXiv:1010.3967v2}.

\bibitem{KapranovChow}
Mikhail M. Kapranov,
\emph{Chow quotients of Grassmannians. I},
in \emph{I.~M.~Gelfand Seminar},
Advances in Soviet Mathematics, vol.~16, part~2,
American Mathematical Society, Providence, RI, 1993, pp.~29--110.
\href{https://arxiv.org/abs/alg-geom/9210002}{arXiv:alg-geom/9210002}.

\bibitem{Aksnes}
E. Aksnes,
\emph{Tropicalization of curve arrangement complements and arroids},
Mathematica Scandinavica \textbf{131} (2025), no.~2.
\url{https://doi.org/10.7146/math.scand.a-156674};
\url{https://arxiv.org/abs/2404.14380v3}.

\bibitem{MichalakReeb}
\L{}. P. Michalak,
\emph{Realization of a graph as the Reeb graph of a Morse function on a manifold},
Topological Methods in Nonlinear Analysis \textbf{52} (2018), no.~2, 749--762.
Definition~2.2 and Theorem~5.6.
\url{https://doi.org/10.12775/TMNA.2018.029};
\url{https://arxiv.org/abs/1805.06727v2}.

% Add this entry inside the main thebibliography environment.
\bibitem{Sugawara}
S.~Sugawara,
\emph{Handle decompositions and Kirby diagrams for the complements of plane algebraic curves},
arXiv:2306.10519v3, 23 April 2026, 23 pp.
\url{https://arxiv.org/abs/2306.10519v3}.
Generic projection setup: \S2; handle counts: Proposition~3.1 and Theorem~3.7;
framed Kirby algorithm: Theorem~4.1.

\bibitem{Hattori}
A. Hattori,
\emph{Topology of $\mathbb C^n$ minus a finite number of affine hyperplanes in general position},
Journal of the Faculty of Science, University of Tokyo, Section~IA, Mathematics
\textbf{22} (1975), no.~2, 205--219.
\url{https://repository.dl.itc.u-tokyo.ac.jp/records/39761}.


\bibitem{PapadimaSuciu}
S. Papadima and A. I. Suciu,
\emph{Higher homotopy groups of complements of complex hyperplane arrangements},
Advances in Mathematics \textbf{165} (2002), 71--100.
\S1.1 (Hattori's theorem for generic arrangements).
\url{https://arxiv.org/abs/math/0002251}.


\bibitem{AllermannRau}
L. Allermann and J. Rau,
\emph{First steps in tropical intersection theory},
Mathematische Zeitschrift \textbf{264} (2010), no.~3, 633--670.
Definition~3.4, Remark~3.5, Proposition~3.7;
Definition~6.1 and Construction~6.4.
\url{https://doi.org/10.1007/s00209-009-0483-1};
\url{https://arxiv.org/abs/0709.3705v3}.

\bibitem{SymingtonFour}
M. Symington,
\emph{Four dimensions from two in symplectic topology},
arXiv:math/0210033v1 (2002).
Proposition~2.15, \S4.3, equation~(4.10), and Proposition~4.14.
\url{https://arxiv.org/abs/math/0210033v1}.

\bibitem{SergeevSchneiderButkovic}
S. Sergeev, H. Schneider and P. Butkovi\v c,
\emph{On visualization scaling, subeigenvectors and Kleene stars in max algebra},
Linear Algebra Appl. \textbf{431} (2009), no.~12, 2395--2406.
\url{https://doi.org/10.1016/j.laa.2009.03.040};
\url{https://arxiv.org/abs/0808.1992}.
% Theorem 2.2 and Proposition 3.4 checked in arXiv:0808.1992v3.

\bibitem{CorralesValencia}
H. Corrales and C. E. Valencia,
\emph{Arithmetical structures on graphs},
Linear Algebra Appl. \textbf{536} (2018), 120--151.
\url{https://doi.org/10.1016/j.laa.2017.09.018};
\url{https://arxiv.org/abs/1604.02502}.
% Theorems 3.2 and 3.7 checked in arXiv:1604.02502v3.
% Published numbering differs at the finiteness result; cite section or version.



\bibitem{LinSturmfels}
S. Lin and B. Sturmfels,
\emph{Polynomial relations among principal minors of a $4\times4$-matrix},
J. Algebra \textbf{322} (2009), 4121--4131.
\href{https://doi.org/10.1016/j.jalgebra.2009.06.026}{doi:10.1016/j.jalgebra.2009.06.026}.
\href{https://arxiv.org/abs/0812.0601}{arXiv:0812.0601v3}.


\bibitem{EvansUrzua}
J. D. Evans and G. Urz\'ua,
\emph{Antiflips, mutations, and unbounded symplectic embeddings of rational homology balls},
Ann. Inst. Fourier \textbf{71} (2021), no.~5, 1807--1843.
\url{https://doi.org/10.5802/aif.3429}.

\bibitem{BHS}
J. Brendel, J. Hauber and J. Schmitz,
\emph{Semi-Local Exotic Lagrangian Tori in Dimension Four},
Ann. Inst. Fourier, online first (2026), 48 pp.
\url{https://doi.org/10.5802/aif.3789};
\url{https://arxiv.org/abs/2403.00408}.

\bibitem{RuddatAnnouncement2026}
H. Ruddat,
\emph{Lagrangian torus fibrations of Calabi--Yau hypersurfaces},
Institut Mittag--Leffler, seminar abstract, 30 June 2026.
Joint work with C. Y. Mak, D. Matessi and I. Zharkov, as stated in the abstract.
Announcement; accessed 28 October 2026.
\url{https://www.mittag-leffler.se/seminar/helge-ruddat-tba/}.

\bibitem{ZharkovSoftCells2026}
I. Zharkov,
\emph{Phase Tropical hypersurfaces and their cousins},
ICMS--Sofia Geometry Seminar, 15 April 2026; announcement posted 8 April 2026.
Seminar abstract; accessed 28 October 2026.
\href{https://icms.bg/phase-tropical-hypersurfaces-and-their-cousins-geometry-seminar-talk-by-ilia-zharkov/}{ICMS--Sofia: seminar announcement}.


\bibitem{VenugopalanWoodward2026}
S. Venugopalan and C. T. Woodward,
\emph{Tropical disk potential for almost toric manifolds},
arXiv:2604.03161v1, 3 April 2026.
Definitions~1.1--1.2, Theorem~1.3, and the discussion of the nonmonotone potential in \S1.
\url{https://arxiv.org/abs/2604.03161v1}.

\end{thebibliography}

\end{document}

